\documentclass[twoside,11pt]{article}

\usepackage{blindtext}

% Any additional packages needed should be included after jmlr2e.
% Note that jmlr2e.sty includes epsfig, amssymb, natbib and graphicx,
% and defines many common macros, such as 'proof' and 'example'.
%
% It also sets the bibliographystyle to plainnat; for more information on
% natbib citation styles, see the natbib documentation, a copy of which
% is archived at http://www.jmlr.org/format/natbib.pdf

% Available options for package jmlr2e are:
%
%   - abbrvbib : use abbrvnat for the bibliography style
%   - nohyperref : do not load the hyperref package
%   - preprint : remove JMLR specific information from the template,
%         useful for example for posting to preprint servers.
%
% Example of using the package with custom options:
%
% \usepackage[abbrvbib, preprint]{jmlr2e}

\usepackage{jmlr2e}

% Definitions of handy macros can go here

\newcommand{\dataset}{{\cal D}}
\newcommand{\fracpartial}[2]{\frac{\partial #1}{\partial  #2}}

% Heading arguments are {volume}{year}{pages}{date submitted}{date published}{paper id}{author-full-names}

\usepackage{lastpage}
\jmlrheading{23}{2022}{1-\pageref{LastPage}}{1/21; Revised 5/22}{9/22}{21-0000}{Author One and Author Two}

% Short headings should be running head and authors last names
\ShortHeadings{Sample JMLR Paper}{One and Two}
\firstpageno{1}

% custom pacages and commands


% 1. Mathematical Packages
\usepackage{amsmath, mathtools}
\usepackage{bm}

% 2. Algorithms
\usepackage{algorithm}
\usepackage{algpseudocode}

\algrenewcommand\algorithmicrequire{\textbf{Require:}}
\algrenewcommand\algorithmicensure{\textbf{Ensure:}}

% 3. Figures and Tables
\usepackage{subcaption}
\usepackage{booktabs}

% 42. external file
\input{mathdef.tex}

\begin{document}

\title{formulation memo}

\author{\name Ryogo Tanaka \email @osaka-u \\
       \addr Dep.\\
       Habdai.\\
       Suita.
       % \AND
       % \name Author Two \email two@cs.berkeley.edu \\
       % \addr Division of Computer Science\\
       % University of California\\
       % Berkeley, CA 94720-1776, USA
       }

\editor{My editor}

\maketitle

% \begin{abstract}%   <- trailing '%' for backward compatibility of .sty file
% \end{abstract}

% \begin{keywords}
%   keyword one, keyword two, keyword three
% \end{keywords}

\section{Motivation and positioning}
\label{sec:motivation}

Learning stochastic nonlinear dynamics from high-dimensional and
partially observed time series requires both a representation of the
latent dynamical state and an operator that propagates that
representation forward in time. Embedded latent transfer-operator
models address this problem by representing latent state laws through
Hilbert-space mean embeddings and by propagating these embedded laws
with linear transfer operators. In this view, the object being propagated
is not a point-valued latent state, but an embedded probability law.

Deep Spectral Learning of Embedded Latent Transfer Operators (DSE)
extends this operator-theoretic viewpoint to learned deep features and
finite-dimensional Galerkin coordinates. Given a learned observation
encoder, DSE constructs predictive state coordinates by a functional
CCA / stochastic realization step and then estimates the embedded
transfer and observable operators by ridge-regularized Galerkin
regressions. This is a powerful spectral construction: the latent
coordinates are chosen so that past and future observation features are
maximally correlated in an appropriate feature space.

The goal of proposed method is not to reject this spectral state construction.
Rather, we seek an alternative state-construction principle whose
semantics are tied directly to filtering. In DSE, the CCA-based
coordinates are constructed before the filtering problem is explicitly
formulated. Consequently, the training-time state construction and the
test-time prediction-correction update are conceptually separated. In
contrast, we would like the latent coordinates themselves to be defined
as prior and posterior embedded beliefs.

To this end, we introduce two Galerkin coordinates at each time index:
\begin{align}
\va_t^-
&\approx
P_A
\mathbb E[
\phi_x(\rvx(t))
\mid
\rvy(1{:}t-1)
],
\label{eq:motivation_prior_coordinate}
\\
\va_t^+
&\approx
P_A
\mathbb E[
\phi_x(\rvx(t))
\mid
\rvy(1{:}t)
].
\label{eq:motivation_posterior_coordinate}
\end{align}
Here, $\va_t^-$ is intended to represent the projected prior embedded
state law before incorporating the observation at time $t$, while
$\va_t^+$ is intended to represent the projected posterior embedded
state law after incorporating that observation. Thus, the state
coordinates are not point-valued latent states. They are relaxed
Galerkin coordinates of projected conditional mean embeddings.

This perspective is inspired by kernel Bayes and kernel Kalman
inference, where probability laws are represented by RKHS mean
embeddings and belief updates are performed through Hilbert-space
residual minimization rather than explicit parametric density
estimation. We do not aim to implement the kernel Kalman rule directly.
Instead, we use its prior/posterior embedded-belief update philosophy as
a state-construction principle within the Galerkin operator-learning
framework of DSE.

The proposed formulation replaces the CCA-based realization step by an
embedding free energy over the prior and posterior coordinates
$\{\va_t^-,\va_t^+\}_{t=1}^T$. This free energy consists of three
components: observation consistency, prior-to-posterior correction, and
transfer consistency. Observation consistency requires the posterior
coordinate to explain the observed embedding through the observable
operator. Prior-to-posterior correction prevents the posterior embedded
belief from moving unnecessarily far from the prior embedded belief.
Transfer consistency requires the posterior coordinate at time $t$ to be
propagated by the learned transfer operator to the prior coordinate at
time $t+1$.

The transfer and observable matrices are still estimated as Galerkin
operators. In particular, $V_A$ and $V_B$ are not arbitrary neural
layers; they are ridge-regularized finite-dimensional representations of
projected conditional operators. Thus, the operator-learning core of DSE
is retained, while the CCA-based state construction is replaced by a
Hilbert-space embedded-belief inference principle.

After learning the embedded-state coordinates and operators, prediction
is performed by a causal filtering recursion. The learned transfer
operator propagates the posterior coordinate from the previous time step
to a new prior coordinate,
\begin{align}
\bar{\va}_t^-
=
V_A \bar{\va}_{t-1}^+ .
\label{eq:motivation_test_prior}
\end{align}
The current observation then corrects this prior through an
embedding-space innovation. In the centered Galerkin convention used
below, this correction has the form
\begin{align}
\bar{\va}_t^+
=
\bar{\va}_t^-
+
K_t
\left(
\tilde{\vb}_t - V_B\bar{\va}_t^-
\right),
\label{eq:motivation_test_posterior}
\end{align}
where $\tilde{\vb}_t$ denotes the centered observation-embedding
coordinate. The prior coordinate $\bar{\va}_t^-$, not the posterior
coordinate, is used for one-step causal prediction through the
observable operator and a read-out / decoder map. This separates
embedded-state and operator learning from raw-observation prediction,
while preserving the same prediction--correction semantics at training
and test time.

The resulting method is therefore not a new non-Gaussian filtering
method in isolation, nor a direct deep implementation of KKR. It is a
Galerkin embedded-belief learning formulation that connects the
Hilbert-space filtering philosophy of kernel Kalman methods with the
operator-learning structure of DSE. Its central aim is to make state
construction, observation correction, transfer consistency, and
test-time causal prediction part of a single operator-theoretic
formulation.


\section{Background overview}
\label{sec:background}

A probabilistic state-space model is specified by a latent process
$\{\rvx(t)\}_{t=1}^T$ and an observation process
$\{\rvy(t)\}_{t=1}^T$, with realizations
$\vx(t)\in\mathcal X$ and $\vy(t)\in\mathcal Y$. Its joint law is written as
\begin{align}
p_\theta(\rvx(1{:}T),\rvy(1{:}T))
=
p_\theta(\rvx(1))
\prod_{t=2}^T p_\theta(\rvx(t)\mid \rvx(t-1))
\prod_{t=1}^T p_\theta(\rvy(t)\mid \rvx(t)).
\label{eq:ssm_factorization}
\end{align}
This factorization separates temporal propagation through the transition law from data generation through the observation law. In the linear-Gaussian case,
\begin{align}
\rvx(t)
&=
\mA \rvx(t-1) + \rvw(t),
&
\rvy(t)
&=
\mC \rvx(t) + \rvv(t),
\label{eq:linear_gaussian_ssm}
\end{align}
with $\rvw(t)\sim\mathcal N(0,\mQ)$ and $\rvv(t)\sim\mathcal N(0,\mR)$, the complete-data negative log density is, up to constants,
\begin{align}
-\log p_\theta(\rvx(1{:}T),\rvy(1{:}T))
&=
\frac{1}{2}
\sum_{t=1}^T
(\vy(t)-\mC \vx(t))^\top \mR^{-1}(\vy(t)-\mC \vx(t))
\nonumber\\
&\quad+
\frac{1}{2}
\sum_{t=2}^T
(\vx(t)-\mA \vx(t-1))^\top \mQ^{-1}(\vx(t)-\mA \vx(t-1))
\nonumber\\
&\quad+
\frac{1}{2}
(\vx(1)-\bm{\pi}_1)^\top \mV_1^{-1}(\vx(1)-\bm{\pi}_1).
\label{eq:complete_data_nll}
\end{align}
The observation and transition terms in \eqref{eq:complete_data_nll} are the quadratic residuals underlying Kalman filtering, smoothing, and maximum-likelihood estimation in linear dynamical systems \citep{katayama2005subspace}.

A complementary description is obtained in Hilbert feature spaces. Let $\rvx$ and $\rvy$ be random variables taking values $\vx \in \mathcal X$ and $\vy \in \mathcal Y$, and let $\mathcal H$ and $\mathcal G$ be separable Hilbert spaces. Consider feature maps
\begin{align}
\phi_x:\mathcal X\to\mathcal H,
\qquad
\psi_y:\mathcal Y\to\mathcal G
\label{eq:feature_maps}
\end{align}
such that
\begin{align}
\mathbb E[\|\phi_x(\rvx)\|_{\mathcal H}^2] < \infty,
\qquad
\mathbb E[\|\psi_y(\rvy)\|_{\mathcal G}^2] < \infty.
\label{eq:feature_square_integrable}
\end{align}
The corresponding mean elements are
\begin{align}
\mu_{\rvx}
&:=
\mathbb E[\phi_x(\rvx)]
=
\int \phi_x(\vx)\,dP_{\rvx}(\vx)
\in\mathcal H,
\nonumber\\
\mu_{\rvy}
&:=
\mathbb E[\psi_y(\rvy)]
=
\int \psi_y(\vy)\,dP_{\rvy}(\vy)
\in\mathcal G.
\label{eq:mean_elements}
\end{align}
For $a\in\mathcal G$ and $b,w\in\mathcal H$, define the rank-one operator
\begin{align}
(a\otimes b)w
:=
\langle w,b\rangle_{\mathcal H}\,a.
\label{eq:rank_one_operator}
\end{align}
The covariance and cross-covariance operators are
\begin{align}
\mathcal C_{\rvx}
&:=
\mathbb E[\phi_x(\rvx)\otimes \phi_x(\rvx)],
&
\mathcal C_{\rvy\rvx}
&:=
\mathbb E[\psi_y(\rvy)\otimes \phi_x(\rvx)].
\label{eq:covariance_crosscovariance}
\end{align}
Restricting conditional expectation to linear predictors on feature space yields the regularized least-squares problem
\begin{align}
\mathcal A^\ast
:=
\arg\min_{\mathcal A\in L(\mathcal H,\mathcal G)}
\mathbb E
\!\left[
\|\psi_y(\rvy)-\mathcal A \phi_x(\rvx)\|_{\mathcal G}^2
\right]
+
\lambda \|\mathcal A\|_{\mathrm{HS}}^2,
\label{eq:conditional_operator_regression}
\end{align}
where $\lambda>0$ and $L(\mathcal H,\mathcal G)$ denotes the space of bounded linear operators from $\mathcal H$ to $\mathcal G$. The normal equation gives
\begin{align}
\mathcal A^\ast
=
\mathcal C_{\rvy\rvx}
(\mathcal C_{\rvx}+\lambda \mathcal I)^{-1},
\label{eq:conditional_operator_solution}
\end{align}
where $\mathcal I$ is the identity operator on $\mathcal H$. Accordingly, the regularized conditional covariance operator is
\begin{align}
\mathcal C_{\rvy\mid \rvx}
:=
\mathcal C_{\rvy\rvx}
(\mathcal C_{\rvx}+\lambda \mathcal I)^{-1}.
\label{eq:conditional_covariance_operator}
\end{align}
This operator is the population-optimal linear predictor of $\psi_y(\rvy)$ from $\phi_x(\rvx)$ in the regularized mean-squared sense \citep{Fukumizu2004,Hsing2015,Song2009}.

The operator-valued state-space representation used in DSE is obtained by applying this construction to stochastic dynamics. Let $\{\rvx(t)\}_{t\in\mathbb T}$ be a latent stationary Markov process on $\mathcal X$, and let $\{\rvy(t)\}_{t\in\mathbb T}$ be an observation process on $\mathcal Y$. At the distribution level, the transition and observation channels map marginal laws by
\begin{align}
p_{t+1}(\vx)
&=
\int_{\mathcal X}
p_{\mathrm{tr}}(\vx \mid \vz)\,p_t(\vz)\,d\vz,
&
q_t(\vy)
&=
\int_{\mathcal X}
p_{\mathrm{ob}}(\vy \mid \vx)\,p_t(\vx)\,d\vx .
\label{eq:density_level_dynamics}
\end{align}
At the embedding level, define
\begin{align}
\mathcal T_e
&:=
\mathcal C_{\rvx(t+1)\mid \rvx(t)},
&
\mathcal O_e
&:=
\mathcal C_{\rvy(t)\mid \rvx(t)}.
\label{eq:operator_level_maps}
\end{align}
Then the embedded laws evolve linearly as
\begin{align}
\mu_{\rvx(t+1)}
&=
\mathcal T_e \mu_{\rvx(t)},
&
\mu_{\rvy(t)}
&=
\mathcal O_e \mu_{\rvx(t)}.
\label{eq:operator_level_ssm}
\end{align}
This is the operator-level state-space model retained in the formulation developed below when stochastic realization is replaced by projected latent-state learning.

At the finite-dimensional level, DSE estimates linear operators on learned feature subspaces by a two-stage deep-feature procedure. For generic random variables $\rvx$ and $\rvz$ and a scalar target $s\in\mathbb R$, the structural objective is
\begin{align}
\min_{\xi,\theta_X,\theta_Z}
\mathbb E
\!\left[
\bigl(
s-r_\xi(\mathbb E[\psi_{\theta_X}(\rvx)\mid \rvz])
\bigr)^2
\right].
\label{eq:two_stage_target}
\end{align}
Its first stage estimates the conditional mean feature map by
\begin{align}
\min_{\mV,\theta_Z}
\frac{1}{N}
\sum_{i=1}^N
\|\psi_{\theta_X}(\vx_i)-\mV\phi_{\theta_Z}(\vz_i)\|_2^2
+
\lambda_1\|\mV\|_F^2,
\label{eq:two_stage_stage1}
\end{align}
whose closed-form solution for fixed features is
\begin{align}
\hat{\mV}
=
\Psi\Phi^\top(\Phi\Phi^\top+\lambda_1\mI)^{-1}.
\label{eq:two_stage_closed_form}
\end{align}
This is the finite-dimensional ridge-regression block that later yields the projected transfer and observable operators. The formulation developed below preserves this operator-estimation mechanism and replaces only the stochastic-realization block by projected latent-state learning.


\section{Statistical Target and Embedding-Free-Energy Principle}
\label{sec:statistical_objective}

The object of inference is not a point-valued latent state, but the
prior and posterior mean embeddings of the latent state law. Let
$\{\rvx(t)\}_{t=1}^T$ be the latent Markov process and
$\{\rvy(t)\}_{t=1}^T$ the observation process. Let
$\phi_x:\mathcal X\to\mathcal H$ and
$\psi_y:\mathcal Y\to\mathcal G$ be square-integrable feature maps.
For each time index $t$, define the ideal prior and posterior embedded
state laws by
\begin{align}
\mu_{x,t}^{-,\star}
&:=
\mathbb E[
\phi_x(\rvx(t))
\mid
\rvy(1{:}t-1)
],
\label{eq:ideal_prior_embedding}
\\
\mu_{x,t}^{+,\star}
&:=
\mathbb E[
\phi_x(\rvx(t))
\mid
\rvy(1{:}t)
].
\label{eq:ideal_posterior_embedding}
\end{align}
These are Hilbert-space valued conditional expectations. Equivalently,
they are the orthogonal projections of $\phi_x(\rvx(t))$ onto the
closed subspaces of Hilbert-space valued random variables measurable
with respect to $\sigma(\rvy(1{:}t-1))$ and $\sigma(\rvy(1{:}t))$,
respectively. Thus, the statistical target is a conditional mean
embedding, not a realized latent coordinate.

Let $\mathcal H_A\subset\mathcal H$ and
$\mathcal G_B\subset\mathcal G$ be finite-dimensional trial spaces, and
let
\begin{align}
P_A:\mathcal H\to\mathcal H_A,
\qquad
P_B:\mathcal G\to\mathcal G_B
\label{eq:projection_operators_statistical}
\end{align}
be the corresponding orthogonal projections. The finite-dimensional
targets are the projected prior and posterior mean embeddings
\begin{align}
\bar\mu_{t}^{-,\star}
&:=
P_A\mu_{x,t}^{-,\star},
&
\bar\mu_{t}^{+,\star}
&:=
P_A\mu_{x,t}^{+,\star}.
\label{eq:projected_prior_posterior_embeddings}
\end{align}
The optimization variables
\begin{align}
\mu_t^- \in \mathcal H_A,
\qquad
\mu_t^+ \in \mathcal H_A
\label{eq:prior_posterior_variables_hilbert}
\end{align}
are relaxed Galerkin representatives intended to approximate
$\bar\mu_t^{-,\star}$ and $\bar\mu_t^{+,\star}$. We do not require
these variables to be exact projected mean embeddings of valid
probability laws; instead, inference is performed in the relaxed trial
space $\mathcal H_A$.

These variables are training-time variational representatives. In the
core learning problem, the entire sequence
$\{\mu_t^-,\mu_t^+\}_{t=1}^T$ is optimized jointly together with the
projected operators. Therefore, although the variables are assigned
prior/posterior embedded-belief semantics, they should not be identified
with the output of an online filtering recursion.

After the projected operators and the Galerkin coordinate system have
been learned, the causal filtering path is generated separately by a
recursive prediction--correction rule. We denote that path by
$\{\bar\mu_t^-,\bar\mu_t^+\}_{t=1}^T$ in Hilbert-space notation and by
$\{\bar{\va}_t^-,\bar{\va}_t^+\}_{t=1}^T$ in finite-dimensional
coordinates. This distinction separates the batch construction of
embedded-belief coordinates from the causal filtering procedure used for
prediction.

The observation at time $t$ gives the empirical observation-law
embedding
\begin{align}
\psi_y(\vy(t))\in\mathcal G.
\end{align}
Its projected observation feature is
\begin{align}
\psi_t^\Pi
:=
P_B\psi_y(\vy(t))
\in\mathcal G_B.
\label{eq:projected_observation_embedding_revised}
\end{align}
Let
\begin{align}
\mathcal T_A:\mathcal H_A\to\mathcal H_A,
\qquad
\mathcal O_B:\mathcal H_A\to\mathcal G_B
\label{eq:projected_transfer_observable_hilbert}
\end{align}
denote the projected transfer and observable operators. The operator
$\mathcal T_A$ propagates posterior embedded state laws to prior
embedded state laws, while $\mathcal O_B$ maps an embedded state law to
a projected observation-law embedding.

The prior prediction is defined by the projected transfer operator,
\begin{align}
\mu_{t+1}^-
\approx
\mathcal T_A \mu_t^+.
\label{eq:prior_prediction_hilbert}
\end{align}
At time $t$, given a prior embedded state $\mu_t^-$, the posterior
embedded state is defined as the solution of a Hilbert-space
regularized inverse problem. In the ideal noiseless case, one would seek
an element $\mu\in\mathcal H_A$ satisfying the observation consistency
condition
\begin{align}
\mathcal O_B\mu
=
\psi_t^\Pi.
\label{eq:hard_observation_consistency}
\end{align}
Among all such elements, the posterior should remain closest to the
prior embedded state $\mu_t^-$. This gives the constrained projection
problem
\begin{align}
\min_{\mu\in\mathcal H_A}
\frac12
\|\mu-\mu_t^-\|_{(\mathcal C_t^-)^\dagger}^2
\quad
\text{subject to}
\quad
\mathcal O_B\mu=\psi_t^\Pi,
\label{eq:hard_projection_posterior_update}
\end{align}
where $\mathcal C_t^-:\mathcal H_A\to\mathcal H_A$ is a positive
semi-definite metric or preconditioning operator associated with the
prior embedded state.

In practice, the hard constraint in
\eqref{eq:hard_projection_posterior_update} is too restrictive because
of stochastic observations, finite-dimensional projection error, and
operator-estimation error. We therefore use its Tikhonov relaxation:
\begin{align}
\mu_t^+
=
\arg\min_{\mu\in\mathcal H_A}
\frac12
\|\psi_t^\Pi-\mathcal O_B\mu\|_{\mathcal R_B^{-1}}^2
+
\frac12
\|\mu-\mu_t^-\|_{(\mathcal C_t^-)^\dagger}^2
+
\Omega_A(\mu).
\label{eq:hilbert_posterior_update_objective}
\end{align}
Here $\mathcal R_B:\mathcal G_B\to\mathcal G_B$ is a positive definite
metric for observation-embedding residuals, and
$\Omega_A$ is an optional regularizer on the Galerkin state element.
The operator $\mathcal C_t^-$ is not assumed to be the covariance of a
Gaussian latent variable. It specifies the local geometry of the
prior-to-posterior update in the embedded-state space, playing the role
of a recursive least-squares preconditioner or inverse-curvature metric.

When $\Omega_A=0$, the first-order condition for
\eqref{eq:hilbert_posterior_update_objective} is
\begin{align}
\left(
\mathcal O_B^\ast\mathcal R_B^{-1}\mathcal O_B
+
(\mathcal C_t^-)^\dagger
\right)
\mu_t^+
=
\mathcal O_B^\ast\mathcal R_B^{-1}\psi_t^\Pi
+
(\mathcal C_t^-)^\dagger \mu_t^-.
\label{eq:hilbert_posterior_normal_equation}
\end{align}
Equivalently, when the inverses are well-defined on the relevant
subspaces, this update can be written in innovation form as
\begin{align}
\mu_t^+
&=
\mu_t^-
+
\mathcal K_t
\left(
\psi_t^\Pi-\mathcal O_B\mu_t^-
\right),
\label{eq:hilbert_innovation_update}
\\
\mathcal K_t
&=
\mathcal C_t^- \mathcal O_B^\ast
\left(
\mathcal O_B \mathcal C_t^- \mathcal O_B^\ast
+
\mathcal R_B
\right)^{-1}.
\label{eq:hilbert_gain_operator}
\end{align}
This is the Hilbert-space analogue of a Kalman or KKR-type correction:
the update is obtained from a residual-minimization principle on
embedded state laws, not from a Gaussian distributional assumption on
finite-dimensional latent coordinates.

The transfer relation is also imposed at the level of embedded laws. In
an exact infinite-dimensional model one would have
$\mu_{t+1}^-=\mathcal T_A\mu_t^+$. In the projected finite-dimensional
setting, this identity is softened to a transfer-consistency penalty
\begin{align}
\frac12
\|\mu_{t+1}^- - \mathcal T_A\mu_t^+\|_{\mathcal Q_A^{-1}}^2,
\label{eq:hilbert_transfer_consistency}
\end{align}
where $\mathcal Q_A:\mathcal H_A\to\mathcal H_A$ is a positive definite
metric for transfer residuals. This metric represents Galerkin
projection error, operator-estimation error, and model mismatch in the
embedded-state dynamics; it is not required to be interpreted as a
Gaussian process-noise covariance.

For fixed projected operators $\mathcal T_A$ and $\mathcal O_B$, the
training-time embedded-state objective over a sequence is therefore
\begin{align}
\mathcal F_H
(\{\mu_t^-,\mu_t^+\}_{t=1}^T
\mid
\mathcal T_A,\mathcal O_B)
&=
\frac12
\|\mu_1^- - \mu_{\mathrm{init}}\|_{\mathcal P_0^{-1}}^2
\nonumber\\
&\quad+
\sum_{t=1}^T
\left[
\frac12
\|\psi_t^\Pi-\mathcal O_B\mu_t^+\|_{\mathcal R_B^{-1}}^2
+
\frac12
\|\mu_t^+ - \mu_t^-\|_{(\mathcal C_t^-)^\dagger}^2
+
\Omega_A(\mu_t^+)
\right]
\nonumber\\
&\quad+
\sum_{t=1}^{T-1}
\frac12
\|\mu_{t+1}^- - \mathcal T_A\mu_t^+\|_{\mathcal Q_A^{-1}}^2 .
\label{eq:hilbert_embedding_free_energy}
\end{align}
The first line is the initial prior-deviation term. The second line is
the observation-consistency and prior-to-posterior correction energy.
The third line is the transfer-consistency energy. We refer to
\eqref{eq:hilbert_embedding_free_energy} as an embedding free energy:
it has the structure of a variational filtering objective, but its
terms are Hilbert-space discrepancies between embedded laws rather than
negative log likelihoods of an explicit parametric state-space model.

The ridge penalties used to estimate $\mathcal T_A$ and $\mathcal O_B$
are not included in \eqref{eq:hilbert_embedding_free_energy}. They
belong to the separate Galerkin operator-estimation subproblems. This
separation preserves the operator-learning structure used in DSE:
embedded-state inference is performed for fixed projected operators,
and the operators are then updated by regularized conditional-operator
regression.


\section{Finite-Dimensional Galerkin Formulation}
\label{sec:finite_galerkin_formulation}

We now write the embedding free energy
\eqref{eq:hilbert_embedding_free_energy} in finite-dimensional Galerkin
coordinates. Let
\begin{align}
\mathcal H_A
&=
\operatorname{span}\{\chi_1,\dots,\chi_{d_A}\},
&
\mathcal G_B
&=
\operatorname{span}\{\omega_1,\dots,\omega_{d_B}\}
\label{eq:finite_trial_spaces}
\end{align}
be the latent-state and observation trial spaces. For notational
simplicity, we assume in this section that the bases are
orthonormalized. If non-orthonormal bases are used, all coefficient
vectors below are understood as mass-matrix-normalized Galerkin
coordinates.

The projected prior and posterior embedded-state elements are expanded
as
\begin{align}
\mu_t^-
&=
\sum_{i=1}^{d_A}
a_{i,t}^-\,\chi_i,
&
\mu_t^+
&=
\sum_{i=1}^{d_A}
a_{i,t}^+\,\chi_i,
\label{eq:prior_posterior_galerkin_expansion}
\end{align}
with coefficient vectors
\begin{align}
\va_t^-
&=
[a_{1,t}^-,\dots,a_{d_A,t}^-]^\top
\in\mathbb R^{d_A},
&
\va_t^+
&=
[a_{1,t}^+,\dots,a_{d_A,t}^+]^\top
\in\mathbb R^{d_A}.
\label{eq:prior_posterior_coefficient_vectors}
\end{align}
The projected observation embedding is expanded as
\begin{align}
\psi_t^\Pi
=
\sum_{j=1}^{d_B}
b_{j,t}\,\omega_j,
\qquad
\vb_t
=
[b_{1,t},\dots,b_{d_B,t}]^\top
\in\mathbb R^{d_B}.
\label{eq:observation_galerkin_expansion_revised}
\end{align}
The coordinates $\va_t^-$ and $\va_t^+$ are not point-valued latent
states. They are relaxed Galerkin coordinates of projected prior and
posterior conditional mean embeddings. In general, not every vector in
$\mathbb R^{d_A}$ corresponds to the projected mean embedding of a
valid probability law. We therefore optimize over the relaxed Galerkin
space rather than imposing hard probability-simplex or positivity
constraints.

The observation feature coefficients are constructed through a
two-stage observation path. First, a time-invariant encoder maps the raw
observation to an encoder feature,
\begin{align}
\vm_t
=
u_\eta(\vy(t))
\in\mathbb R^m.
\label{eq:encoder_feature_revised}
\end{align}
Second, an observation dictionary maps this encoder feature to the
Galerkin coordinate of the projected observation embedding,
\begin{align}
\vb_t
=
\psi_\omega(\vm_t)
=
\psi_\omega(u_\eta(\vy(t)))
\in\mathbb R^{d_B}.
\label{eq:observation_dictionary_revised}
\end{align}
The coordinate $\vb_t$ represents the projected empirical observation
embedding $\psi_t^\Pi=P_B\psi_y(\vy(t))$ in the chosen observation
trial space. The two-stage path is retained because the operator
learning problem is posed in the observation-embedding coordinates
$\vb_t$, whereas the final prediction is decoded through the encoder
feature $\vm_t$.

For a training sequence, define
\begin{align}
M
&:=
[\vm_1,\dots,\vm_T]
\in\mathbb R^{m\times T},
&
B
&:=
[\vb_1,\dots,\vb_T]
\in\mathbb R^{d_B\times T}.
\label{eq:M_B_observation_matrices}
\end{align}
We use a centered Galerkin convention on the observation side. Let
\begin{align}
\bar{\vm}
&:=
\frac1T\sum_{t=1}^T\vm_t,
&
\tilde{\vm}_t
&:=
\vm_t-\bar{\vm},
&
\tilde M
&:=
[\tilde{\vm}_1,\dots,\tilde{\vm}_T],
\label{eq:centered_encoder_feature_matrix}
\\
\bar{\vb}
&:=
\frac1T\sum_{t=1}^T\vb_t,
&
\tilde{\vb}_t
&:=
\vb_t-\bar{\vb},
&
\tilde B
&:=
[\tilde{\vb}_1,\dots,\tilde{\vb}_T].
\label{eq:centered_observation_embedding_matrix}
\end{align}
The core embedded-state and operator-learning objective is written in
terms of the centered observation embeddings $\tilde{\vb}_t$. The means
$\bar{\vb}$ and $\bar{\vm}$ are stored: $\bar{\vb}$ is used to center
new observation embeddings at test time, and $\bar{\vm}$ is restored
before applying the decoder. This centering convention fixes a numerical
Galerkin gauge for a linear observable operator without introducing an
additional probabilistic assumption on the latent process.

Let
\begin{align}
V_A
&\in\mathbb R^{d_A\times d_A},
&
V_B
&\in\mathbb R^{d_B\times d_A}
\label{eq:finite_operator_matrices_revised}
\end{align}
be the coordinate matrices of the projected transfer and observable
operators,
\begin{align}
V_A
&\leftrightarrow
\mathcal T_A:\mathcal H_A\to\mathcal H_A,
&
V_B
&\leftrightarrow
\mathcal O_B:\mathcal H_A\to\mathcal G_B.
\label{eq:operator_matrix_correspondence}
\end{align}
The transfer matrix maps a posterior embedded-state coordinate at time
$t$ to a prior embedded-state coordinate at time $t+1$, while the
observable matrix maps an embedded-state coordinate to a projected
observation embedding.

For the latent embedded-state coordinates, define
\begin{align}
A^-
&:=
[\va_1^-,\dots,\va_T^-]
\in\mathbb R^{d_A\times T},
&
A^+
&:=
[\va_1^+,\dots,\va_T^+]
\in\mathbb R^{d_A\times T}.
\label{eq:A_minus_A_plus_matrices}
\end{align}
The training-time latent variables are the two coefficient sequences
$A^-$ and $A^+$. The former represents prior or predictive embedded
beliefs, and the latter represents posterior or corrected embedded
beliefs. They are optimized as batch variational coordinates. The
causal filtering coordinates used for prediction are generated later by
the learned recursion and are denoted by $\bar A^-$ and $\bar A^+$.

Let
\begin{align}
R_B
&\in\mathbb S_{++}^{d_B},
&
Q_A
&\in\mathbb S_{++}^{d_A},
&
P_t^-
&\in\mathbb S_{+}^{d_A},
&
P_0
&\in\mathbb S_{++}^{d_A}
\label{eq:finite_metric_matrices}
\end{align}
be the finite-dimensional metrics corresponding to
$\mathcal R_B$, $\mathcal Q_A$, $\mathcal C_t^-$, and
$\mathcal P_0$. We use the notation
\begin{align}
\|x\|_{M}^{2}
:=
x^\top M x.
\label{eq:finite_weighted_norm}
\end{align}
If $P_t^-$ is singular or low rank, $(P_t^-)^\dagger$ denotes its
Moore--Penrose pseudoinverse.

For fixed $V_A$ and $V_B$, the finite-dimensional embedding free energy
under the centered observation convention is
\begin{align}
\mathcal F_{\mathrm{emb}}
(A^-,A^+\mid V_A,V_B)
&=
\frac12
\|\va_1^- - \va_{\mathrm{init}}\|_{P_0^{-1}}^2
\nonumber\\
&\quad+
\sum_{t=1}^{T}
\left[
\frac12
\|\tilde{\vb}_t - V_B\va_t^+\|_{R_B^{-1}}^2
+
\frac12
\|\va_t^+ - \va_t^-\|_{(P_t^-)^\dagger}^2
+
\Omega_A(\va_t^+)
\right]
\nonumber\\
&\quad+
\sum_{t=1}^{T-1}
\frac12
\|\va_{t+1}^- - V_A\va_t^+\|_{Q_A^{-1}}^2 .
\label{eq:finite_embedding_free_energy}
\end{align}
The first term anchors the initial prior embedded-state coordinate. The
second line is the finite-dimensional posterior-correction energy: the
posterior coordinate must explain the centered observation embedding
through $V_B$, while remaining close to the prior coordinate under the
metric $(P_t^-)^\dagger$. The third line enforces transfer consistency
between posterior and next-prior coordinates.

The initial coordinate has a distinct semantic role. Ideally,
$\va_{\mathrm{init}}$ is the Galerkin coordinate of an initial prior
embedded state law,
\begin{align}
\va_{\mathrm{init}}
\approx
P_A\mu_{\mathrm{init}},
\qquad
\mu_{\mathrm{init}}
=
\mathbb E[\phi_x(\rvx(1))\mid \mathcal C_0],
\label{eq:initial_prior_embedding_coordinate}
\end{align}
where $\mathcal C_0$ denotes information available before the first
observation in the sequence. If a meaningful initial embedded law is
available, its coordinate can be supplied directly to the algorithm.

When such an initial embedded law is unavailable, we use the origin of
the centered Galerkin gauge as a reference prior and set
\begin{align}
\va_{\mathrm{init}}=0.
\label{eq:zero_reference_initial_coordinate}
\end{align}
This fallback should not be read as a claim that the true initial state
law is represented by the zero vector. In that case, the initial term is
a weak boundary regularization for the first prior coordinate. Its
strength is controlled by the metric $P_0$.

The optional term $\Omega_A$ is used only when additional coordinate
regularization is desired. In the main formulation, we take
\begin{align}
\Omega_A(\va)=0.
\label{eq:default_omega_A_zero}
\end{align}
This choice leaves the coordinate variables as relaxed Galerkin
representatives of projected mean embeddings. Stronger validity
constraints, such as nonnegative or simplex constraints on atom weights,
are not part of the main method.

The metrics $R_B$, $Q_A$, $P_t^-$, and $P_0$ specify the geometry of
the embedding discrepancies. They are not introduced as covariance
parameters of a linear-Gaussian latent state-space model. In the main
fixed-metric formulation, we use scalar metrics
\begin{align}
R_B
&=
\sigma_B^2 I_{d_B},
&
Q_A
&=
\sigma_A^2 I_{d_A},
&
P_t^-
&=
p I_{d_A},
&
P_0
&=
c_0p I_{d_A}.
\label{eq:fixed_scalar_metrics}
\end{align}
The scalars $\sigma_B^2$, $\sigma_A^2$, $p$, and $c_0$ are
user-specified hyperparameters held fixed during the core optimization.
Equivalently, the first three scalars define the loss weights
\begin{align}
w_B
=
\sigma_B^{-2},
\qquad
w_P
=
p^{-1},
\qquad
w_A
=
\sigma_A^{-2}.
\label{eq:fixed_metric_loss_weights}
\end{align}
The parameter $c_0$ controls the strength of the initial anchor. When
$\va_{\mathrm{init}}$ is supplied from a meaningful initial embedded
law, $c_0$ reflects the confidence assigned to that initial law. When
$\va_{\mathrm{init}}=0$ is used only as a reference prior in the
centered Galerkin gauge, the anchor should be interpreted as weak
boundary regularization.

Residual-calibrated metrics and KKR-style covariance or
inverse-curvature recursions may be introduced as metric variants. They
modify the geometry of the embedding free energy and the correction
gain, but they do not change the statistical target.

The operator matrices $V_A$ and $V_B$ are not treated as arbitrary
neural layers. They are Galerkin representations of projected
conditional operators and are estimated by separate ridge-regularized
operator subproblems. Consequently, the Hilbert--Schmidt or Frobenius
regularization on $V_A$ and $V_B$ is attached to those operator
subproblems, rather than being added as an additional penalty inside
\eqref{eq:finite_embedding_free_energy}. This separation mirrors the
operator-learning structure of DSE while replacing its CCA-based state
construction by the embedded-state inference objective above.





\section{Core Learning Algorithm}
\label{sec:core_learning_algorithm}

This section gives the training-time algorithm for minimizing the
finite-dimensional embedding free energy in
\eqref{eq:finite_embedding_free_energy}. The core solver estimates the
training-time variational coordinate arrays $A^-$ and $A^+$ together
with the projected transfer and observable operators $V_A$ and $V_B$.
The coordinates are optimized in batch and are assigned
prior/posterior embedded-belief semantics through the objective; the
causal filtering coordinates used for prediction are generated later by
the learned recursion.

The encoder $u_\eta$ and the observation dictionary $\psi_\omega$ are
held fixed in this core block. Thus, the core solver operates on a fixed
centered observation-embedding matrix $\tilde B$. Raw-observation
prediction, decoder fitting, and any controlled refinement of the
encoder or dictionary are introduced only after the embedded-state and
operator-learning block has been estimated.

\subsection{Observation Features}

Given observations $\{\vy(t)\}_{t=1}^T$, compute encoder features and
observation-embedding coordinates by
\begin{align}
\vm_t
&=
u_\eta(\vy(t)),
&
\vb_t
&=
\psi_\omega(\vm_t).
\label{eq:core_observation_features}
\end{align}
The two-stage observation path is retained because the core
operator-learning problem is posed in the observation-embedding
coordinates $\vb_t$, whereas the final prediction is decoded through
the encoder-feature coordinates $\vm_t$.

Stack the encoder features and observation embeddings as
\begin{align}
M
&=
[\vm_1,\dots,\vm_T]
\in\mathbb R^{m\times T},
&
B
&=
[\vb_1,\dots,\vb_T]
\in\mathbb R^{d_B\times T}.
\label{eq:core_M_B_matrices}
\end{align}
The core solver uses centered training targets. Define
\begin{align}
\bar{\vm}
&=
\frac1T\sum_{t=1}^T \vm_t,
&
\tilde{\vm}_t
&=
\vm_t-\bar{\vm},
&
\tilde M
&=
[\tilde{\vm}_1,\dots,\tilde{\vm}_T],
\label{eq:core_centered_M_matrix}
\\
\bar{\vb}
&=
\frac1T\sum_{t=1}^T \vb_t,
&
\tilde{\vb}_t
&=
\vb_t-\bar{\vb},
&
\tilde B
&=
[\tilde{\vb}_1,\dots,\tilde{\vb}_T].
\label{eq:core_centered_B_matrix}
\end{align}
The matrix $\tilde B$ is the fixed observation target used in the
embedded-state and operator-learning block. The means $\bar{\vb}$ and
$\bar{\vm}$ are stored. The former is used to center observation
embeddings at test time, and the latter is restored before applying the
decoder in the raw-prediction stage.

\subsection{Initialization of Embedded-State Coordinates}

Since the true latent states $\vx(t)$ are not observed, the
embedded-state coordinate arrays cannot be initialized from latent-state
samples. The arrays $A^{+,(0)}$ and $A^{-,(0)}$ are numerical seeds for
the alternating minimization of the embedding free energy. They are not
statistical estimators of the prior or posterior embedded beliefs, and
they do not define the latent state-coordinate system.

The observation embedding matrix $\tilde B$ is not used to define the
latent coordinates. It enters the learning problem through the
observation-consistency and correction terms in the embedding free
energy. Thus, the prior/posterior embedded-belief interpretation of
$A^-$ and $A^+$ is imposed by the objective and by the learned
prediction--correction recursion, not by the numerical initializer.

The main formulation only requires a nondegenerate numerical seed for
$A^{+,(0)}$ and $A^{-,(0)}$ so that the alternating solver can start.
The concrete choice of this seed is part of the implementation protocol.
It may be randomized, deterministic, or supplied externally, but it
should not be interpreted as a state-construction principle. In
particular, the main formulation does not use DSE CCA coordinates,
stochastic-realization coordinates, or a spectral decomposition of
$\tilde B$ as the definition of the embedded-state coordinates.

\subsection{Metric Choices}

The core objective is defined for general positive definite metrics
$R_B$ and $Q_A$ and positive semi-definite prior metrics $P_t^-$. The
metric $P_0$ controls the boundary anchor for the first prior
coordinate. In the main fixed-metric formulation, we use scalar metrics
\begin{align}
R_B
&=
\sigma_B^2 I_{d_B},
&
Q_A
&=
\sigma_A^2 I_{d_A},
&
P_t^-
&=
p I_{d_A},
&
P_0
&=
c_0p I_{d_A}.
\label{eq:core_fixed_metrics}
\end{align}
The scalars $\sigma_B^2$, $\sigma_A^2$, $p$, and $c_0$ are
user-specified hyperparameters held fixed during the core coordinate
optimization. Equivalently, the first three scalars define the loss
weights
\begin{align}
w_B
=
\sigma_B^{-2},
\qquad
w_P
=
p^{-1},
\qquad
w_A
=
\sigma_A^{-2}.
\label{eq:core_loss_weights}
\end{align}
These quantities specify the geometry and relative weighting of the
embedding discrepancies. They are not learned covariance parameters of
a linear-Gaussian latent state-space model.

The residual-calibrated variant updates $R_B$ and $Q_A$ from empirical
residuals
\begin{align}
\rho_t
&=
\tilde{\vb}_t - V_B\va_t^+,
&
\epsilon_t
&=
\va_{t+1}^- - V_A\va_t^+.
\label{eq:residual_definitions_metric}
\end{align}
For example, a diagonal residual-calibrated variant uses
\begin{align}
R_B
&=
\operatorname{diag}
\left(
\frac1T
\sum_{t=1}^T
\rho_t\rho_t^\top
\right)
+
\varepsilon_R I,
\label{eq:RB_diagonal_calibration}
\\
Q_A
&=
\operatorname{diag}
\left(
\frac1{T-1}
\sum_{t=1}^{T-1}
\epsilon_t\epsilon_t^\top
\right)
+
\varepsilon_Q I.
\label{eq:QA_diagonal_calibration}
\end{align}
Such calibration is performed in an outer loop, and the metrics are held
fixed during each coordinate-minimization step. A KKR-style
metric-recursive variant may also update $P_t^-$ through a covariance or
inverse-curvature recursion. These variants change the metric geometry
of the embedding free energy and the correction gain, but not the
statistical target.

\subsection{Coordinate Updates for Fixed Operators}

For fixed $V_A$, $V_B$, and metrics, the coordinate variables
$A^-$ and $A^+$ are updated by block coordinate minimization of
\eqref{eq:finite_embedding_free_energy}. In the main formulation we take
$\Omega_A=0$.

\paragraph{Posterior-coordinate update.}
For fixed $A^-$, $V_A$, and $V_B$, each posterior coordinate $\va_t^+$
solves a linear system
\begin{align}
H_t^+ \va_t^+
=
h_t^+,
\label{eq:Aplus_update_system}
\end{align}
where
\begin{align}
H_t^+
&=
V_B^\top R_B^{-1}V_B
+
(P_t^-)^\dagger
+
\mathbf 1_{\{t<T\}}
V_A^\top Q_A^{-1}V_A,
\label{eq:Aplus_update_H}
\\
h_t^+
&=
V_B^\top R_B^{-1}\tilde{\vb}_t
+
(P_t^-)^\dagger \va_t^-
+
\mathbf 1_{\{t<T\}}
V_A^\top Q_A^{-1}\va_{t+1}^-.
\label{eq:Aplus_update_h}
\end{align}
Thus,
\begin{align}
\va_t^+
=
(H_t^+)^{-1}h_t^+.
\label{eq:Aplus_update_closed_form}
\end{align}
This update balances centered observation consistency, prior-to-posterior
proximity, and the ability of the posterior coordinate to predict the
next prior coordinate.

\paragraph{Prior-coordinate update.}
For fixed $A^+$ and $V_A$, each prior coordinate $\va_t^-$ is updated
from the terms involving the correction energy and the transfer
consistency. At the first time index, there is no previous posterior
coordinate to propagate, and the initial-prior anchor is used:
\begin{align}
H_1^- \va_1^-
&=
h_1^-,
\label{eq:Aminus_initial_system}
\\
H_1^-
&=
(P_1^-)^\dagger
+
P_0^{-1},
\label{eq:Aminus_initial_H}
\\
h_1^-
&=
(P_1^-)^\dagger \va_1^+
+
P_0^{-1}\va_{\mathrm{init}},
\label{eq:Aminus_initial_h}
\\
\va_1^-
&=
(H_1^-)^{-1}h_1^-.
\label{eq:Aminus_initial_update}
\end{align}
For $t>1$, the prior coordinate is updated by
\begin{align}
H_t^- \va_t^-
=
h_t^-,
\label{eq:Aminus_update_system}
\end{align}
with
\begin{align}
H_t^-
&=
(P_t^-)^\dagger
+
Q_A^{-1},
\label{eq:Aminus_update_H}
\\
h_t^-
&=
(P_t^-)^\dagger \va_t^+
+
Q_A^{-1}V_A\va_{t-1}^+.
\label{eq:Aminus_update_h}
\end{align}
Hence,
\begin{align}
\va_t^-
=
(H_t^-)^{-1}h_t^-,
\qquad
t=2,\dots,T.
\label{eq:Aminus_update_closed_form}
\end{align}
For multiple trajectories, transition terms are evaluated only on valid
within-trajectory transitions. Each trajectory start is treated as an
initial index and uses the same initial-prior convention.

\subsection{Galerkin Operator Updates}

For fixed $A^-$ and $A^+$, the transfer and observable matrices are
updated by separate Galerkin ridge-regression subproblems.

The transfer operator is estimated from the relation
\begin{align}
\va_{t+1}^-
\approx
V_A\va_t^+.
\end{align}
The ridge problem is
\begin{align}
V_A
=
\arg\min_{V\in\mathbb R^{d_A\times d_A}}
\frac12
\sum_{t=1}^{T-1}
\|\va_{t+1}^- - V\va_t^+\|_2^2
+
\frac{\lambda_A}{2}
\|V\|_F^2.
\label{eq:VA_operator_subproblem_revised}
\end{align}
Its closed-form solution is
\begin{align}
V_A
=
A_{2:T}^-
(A_{1:T-1}^+)^\top
\left(
A_{1:T-1}^+
(A_{1:T-1}^+)^\top
+
\lambda_A I
\right)^{-1}.
\label{eq:VA_closed_form_revised}
\end{align}
For multiple trajectories, the matrices in this expression are formed
only from valid within-trajectory transition pairs.

The observable operator is estimated from the centered relation
\begin{align}
\tilde{\vb}_t
\approx
V_B\va_t^+.
\end{align}
The ridge problem is
\begin{align}
V_B
=
\arg\min_{V\in\mathbb R^{d_B\times d_A}}
\frac12
\sum_{t=1}^{T}
\|\tilde{\vb}_t - V\va_t^+\|_2^2
+
\frac{\lambda_B}{2}
\|V\|_F^2.
\label{eq:VB_operator_subproblem_revised}
\end{align}
Its closed-form solution is
\begin{align}
V_B
=
\tilde B(A^+)^\top
\left(
A^+(A^+)^\top
+
\lambda_B I
\right)^{-1}.
\label{eq:VB_closed_form_revised}
\end{align}
The ridge penalties in these subproblems are the operator-regression
regularizers. They are not added again as separate terms in the
embedding free energy. If full non-isotropic metrics are used in the
operator subproblems, the corresponding weighted least-squares systems
must be solved directly or reduced to an equivalent whitened problem.
The main core solver uses scalar or diagonal metric conventions that
preserve the closed-form ridge structure.

\subsection{Core Algorithm}

Algorithm~\ref{alg:core_embedding_learning} summarizes the core learning
procedure. During this stage, $A^-$ and $A^+$ are variational
coordinates. They are not tied by the hard equality
$\va_{t+1}^-=V_A\va_t^+$; instead, this relation is imposed softly
through the transfer-consistency term. This avoids a tautological
operator update. After learning, the same transfer relation is used
causally to propagate filtering coordinates.

\begin{algorithm}[t]
\caption{Core Embedded-State Galerkin Learning}
\label{alg:core_embedding_learning}
\begin{algorithmic}[1]
\Require Observations $\{\vy(t)\}_{t=1}^T$; fixed encoder $u_\eta$;
fixed observation dictionary $\psi_\omega$; coordinate dimensions
$d_A,d_B$; ridge parameters $\lambda_A,\lambda_B$; scalar metric
hyperparameters $\sigma_B^2,\sigma_A^2,p,c_0$; optional initial prior
coordinate $\va_{\mathrm{init}}$; number of outer iterations $K$.
\Ensure Embedded-state coordinates $A^-,A^+$, operators $V_A,V_B$, and
stored means $\bar{\vb},\bar{\vm}$.

\State Compute $\vm_t\gets u_\eta(\vy(t))$ and
$\vb_t\gets \psi_\omega(\vm_t)$ for $t=1,\dots,T$.
\State Compute and store $\bar{\vm}$ and $\bar{\vb}$.
\State Form centered variables
$\tilde{\vm}_t\gets\vm_t-\bar{\vm}$ and
$\tilde{\vb}_t\gets\vb_t-\bar{\vb}$.
\State Form $\tilde B\gets[\tilde{\vb}_1,\dots,\tilde{\vb}_T]$.
\State Set
$R_B\gets\sigma_B^2I_{d_B}$,
$Q_A\gets\sigma_A^2I_{d_A}$,
$P_t^-\gets pI_{d_A}$, and
$P_0\gets c_0pI_{d_A}$.
\If{no meaningful initial embedded law is supplied}
    \State Set $\va_{\mathrm{init}}\gets 0$ in the centered Galerkin gauge.
\EndIf
\State Provide nondegenerate numerical seeds
$A^{+,(0)}$ and $A^{-,(0)}$ for the alternating solver.

\For{$k=0,\dots,K-1$}
    \State Update $V_A^{(k+1)}$ by Galerkin ridge regression from
    $A_{1:T-1}^{+,(k)}$ to $A_{2:T}^{-,(k)}$.
    \State Update $V_B^{(k+1)}$ by Galerkin ridge regression from
    $A^{+,(k)}$ to $\tilde B$.
    \State Update each $\va_t^{+,(k+1)}$ by solving the posterior
    coordinate normal equations.
    \State Update each $\va_t^{-,(k+1)}$ by solving the prior
    coordinate normal equations with the initial-prior anchor at
    $t=1$.
    \State Evaluate and record the embedding free-energy components.
\EndFor

\State \Return $A^{-,(K)}$, $A^{+,(K)}$, $V_A^{(K)}$, $V_B^{(K)}$,
$\bar{\vb}$, and $\bar{\vm}$.
\end{algorithmic}
\end{algorithm}

\section{Test-Time Causal Filtering and Prediction}
\label{sec:test_time_filtering_prediction}

The coordinates $A^-$ and $A^+$ estimated in
Algorithm~\ref{alg:core_embedding_learning} are training-time
variational coordinates. They are optimized jointly with the projected
operators under the embedding free energy. At prediction time, however,
one does not optimize an entire coordinate sequence. Instead, the
learned operators are used recursively to produce causal prior and
posterior embedded-state coordinates.

To distinguish the causal filtering path from the training-time
variational coordinates, we write
\begin{align}
\bar{\va}_t^-,
\qquad
\bar{\va}_t^+
\end{align}
for the causal prior and posterior coordinates. Here
$\bar{\va}_t^-$ is the embedded-state coordinate before incorporating
$\vy(t)$, and $\bar{\va}_t^+$ is the corrected coordinate after the
observation at time $t$ has been incorporated. This path is generated
using the same centered observation convention and the same initial
prior coordinate as in the core formulation.

\subsection{Causal Filtering Recursion}

The filtering recursion is initialized by the initial prior coordinate
\begin{align}
\bar{\va}_1^-
=
\va_{\mathrm{init}}.
\label{eq:test_initial_prior_coordinate}
\end{align}
Ideally, $\va_{\mathrm{init}}$ is supplied as the Galerkin coordinate of
a meaningful initial embedded state law. When such an initial law is not
available, the origin of the centered Galerkin gauge is used as a
reference prior. This fallback is a boundary convention, not a claim
that the true initial state law is represented by the zero vector.

At each time step, the observation is mapped through the fixed encoder
and observation dictionary,
\begin{align}
\vm_t
&=
u_\eta(\vy(t)),
&
\vb_t
&=
\psi_\omega(\vm_t),
\label{eq:test_observation_feature}
\end{align}
and then centered using the training mean stored after feature
preparation:
\begin{align}
\tilde{\vb}_t
=
\vb_t-\bar{\vb}.
\label{eq:test_centered_observation_feature}
\end{align}
The test sequence mean is not used for this centering.

Given the prior coordinate $\bar{\va}_t^-$, the centered observation
prediction and innovation are
\begin{align}
\widehat{\tilde{\vb}}_{t|t-1}
&=
V_B\bar{\va}_t^-,
&
\vr_t
&=
\tilde{\vb}_t
-
V_B\bar{\va}_t^-.
\label{eq:test_innovation}
\end{align}
For a fixed prior metric $P_t^-$ and observation metric $R_B$, the
Galerkin correction gain is
\begin{align}
K_t
=
P_t^- V_B^\top
\left(
V_B P_t^- V_B^\top
+
R_B
\right)^{-1}.
\label{eq:test_gain}
\end{align}
The posterior coordinate is updated by
\begin{align}
\bar{\va}_t^+
=
\bar{\va}_t^-
+
K_t\vr_t.
\label{eq:test_posterior_update}
\end{align}
Equivalently, when $\Omega_A=0$, $\bar{\va}_t^+$ is the minimizer of the
finite-dimensional correction objective
\begin{align}
\bar{\va}_t^+
=
\arg\min_{\va\in\mathbb R^{d_A}}
\frac12
\|\tilde{\vb}_t-V_B\va\|_{R_B^{-1}}^2
+
\frac12
\|\va-\bar{\va}_t^-\|_{(P_t^-)^\dagger}^2 .
\label{eq:test_posterior_argmin}
\end{align}
Thus, the correction step is a finite-dimensional embedded-belief update
driven by a centered observation innovation.

The next prior coordinate is propagated by the learned transfer
operator,
\begin{align}
\bar{\va}_{t+1}^-
=
V_A\bar{\va}_t^+.
\label{eq:test_prior_propagation}
\end{align}
In the main fixed-metric setting, $P_t^-=pI_{d_A}$ and
$R_B=\sigma_B^2I_{d_B}$ are held fixed during the filtering pass. A
metric-recursive variant may instead propagate
\begin{align}
P_t^+
&=
P_t^-
-
K_t V_B P_t^-,
\label{eq:test_cov_update_simple}
\\
P_{t+1}^-
&=
V_A P_t^+ V_A^\top
+
Q_A.
\label{eq:test_cov_prediction}
\end{align}
For numerical stability, the Joseph-form update
\begin{align}
P_t^+
=
(I-K_tV_B)P_t^-(I-K_tV_B)^\top
+
K_tR_BK_t^\top
\label{eq:test_cov_update_joseph}
\end{align}
may be used in such a metric-recursive variant. These metric recursions
are optional refinements of the correction geometry and do not change
the statistical target.

\subsection{Difference Between Training-Time and Test-Time Coordinates}

The distinction between $(A^-,A^+)$ and
$(\bar A^-,\bar A^+)$ is important. During training, the coordinates
$A^-$ and $A^+$ are variational variables in the embedding free-energy
objective. In particular, $A^-$ is not constrained by the hard equality
$\va_{t+1}^-=V_A\va_t^+$. Instead, the transfer relation is imposed
softly by the transfer-consistency term.

At test time, by contrast, the prior coordinate is generated causally by
\eqref{eq:test_prior_propagation}. Thus, the filtered sequence
$\{\bar{\va}_t^-,\bar{\va}_t^+\}$ can be computed online from past
observations and the current observation only. This distinction prevents
the training-time variational path from being confused with a causal
prediction path.

The same convention is used when the causal path is recomputed on the
training sequence: observations are centered using the stored training
mean, and the path is initialized from $\va_{\mathrm{init}}$. Therefore,
evaluation and read-out training are based on the causal path
$\bar A^-,\bar A^+$ rather than on the batch variational coordinates
$A^-,A^+$.

\subsection{Causal Feature Prediction}

After the causal filtering path has been computed, the prior coordinate
$\bar{\va}_t^-$ provides a prediction before $\vy(t)$ is incorporated.
Under the centered observation convention, the predicted observation
embedding is
\begin{align}
\widehat{\tilde{\vb}}_{t|t-1}
=
V_B\bar{\va}_t^-.
\label{eq:predicted_observation_embedding_prior}
\end{align}
A read-out map
\begin{align}
r_{\xi_B}:\mathbb R^{d_B}\to\mathbb R^m
\label{eq:readout_map_definition}
\end{align}
maps this centered operator-projected observation embedding to a
centered encoder-feature prediction:
\begin{align}
\widehat{\tilde{\vm}}_{t|t-1}
=
r_{\xi_B}\!\left(\widehat{\tilde{\vb}}_{t|t-1}\right)
=
r_{\xi_B}(V_B\bar{\va}_t^-).
\label{eq:predicted_centered_encoder_feature}
\end{align}
The feature-level causal prediction loss is
\begin{align}
\mathcal L_{\mathrm{feat}}^{\mathrm{pred}}(\xi_B)
=
\frac1T
\sum_{t=1}^T
\left\|
\tilde{\vm}_t
-
r_{\xi_B}(V_B\bar{\va}_t^-)
\right\|_2^2
+
\lambda_{r,B}\Omega(\xi_B).
\label{eq:feature_prediction_loss}
\end{align}
In the feature read-out stage, the encoder $u_\eta$, observation
dictionary $\psi_\omega$, operators $V_A,V_B$, and causal filtering path
$\{\bar{\va}_t^-,\bar{\va}_t^+\}$ are held fixed, and only
$r_{\xi_B}$ is trained.

If $r_{\xi_B}$ is chosen to be linear, $r_{\xi_B}(h)=R_{\xi_B}h$, the
feature read-out can be estimated by ridge regression. Let
\begin{align}
H^-
&=
[V_B\bar{\va}_1^-,\dots,V_B\bar{\va}_T^-]
\in\mathbb R^{d_B\times T},
&
\tilde M
&=
[\tilde{\vm}_1,\dots,\tilde{\vm}_T]
\in\mathbb R^{m\times T}.
\label{eq:readout_design_matrices}
\end{align}
Then
\begin{align}
R_{\xi_B}
=
\tilde M(H^-)^\top
\left(
H^-(H^-)^\top
+
\lambda_{r,B}I
\right)^{-1}.
\label{eq:linear_readout_solution}
\end{align}
A nonlinear read-out, such as an MLP, may also be used. In that case,
\eqref{eq:feature_prediction_loss} is minimized by gradient descent
while the core operator-learning block remains fixed.

\subsection{Raw Observation Prediction}

After the feature-level read-out has been aligned, the centered feature
prediction is mapped back to the original encoder-feature coordinate
system by restoring the training mean:
\begin{align}
\hat{\vm}_{t|t-1}
=
\bar{\vm}
+
\widehat{\tilde{\vm}}_{t|t-1}
=
\bar{\vm}
+
r_{\xi_B}(V_B\bar{\va}_t^-).
\label{eq:mean_restored_predicted_feature}
\end{align}
A decoder $g_\alpha:\mathbb R^m\to\mathcal Y$ then maps this predicted
encoder feature to the raw observation space:
\begin{align}
\hat{\vy}_{t|t-1}
=
g_\alpha(\hat{\vm}_{t|t-1})
=
g_\alpha\!\left(
\bar{\vm}
+
r_{\xi_B}(V_B\bar{\va}_t^-)
\right).
\label{eq:raw_prediction}
\end{align}
The raw causal prediction loss is
\begin{align}
\mathcal L_{\mathrm{raw}}^{\mathrm{pred}}(\xi_B,\alpha)
=
\frac1T
\sum_{t=1}^T
d_{\mathcal Y}
\left(
\vy(t),
g_\alpha\!\left(
\bar{\vm}
+
r_{\xi_B}(V_B\bar{\va}_t^-)
\right)
\right)
+
\lambda_g\Omega(\alpha)
+
\lambda_{r,B}\Omega(\xi_B),
\label{eq:raw_prediction_loss}
\end{align}
where $d_{\mathcal Y}$ is a task-dependent prediction loss in the
observation space.

The raw-level loss is combined with the centered feature-level anchor:
\begin{align}
\mathcal L_{\mathrm{dec}}^{\mathrm{pred}}
=
\mathcal L_{\mathrm{raw}}^{\mathrm{pred}}
+
\lambda_m
\mathcal L_{\mathrm{feat}}^{\mathrm{pred}}.
\label{eq:decoder_prediction_loss_with_feature_anchor}
\end{align}
The feature anchor keeps the read-out aligned with the encoder-feature
space and prevents the decoder from absorbing arbitrary changes in the
read-out representation.

\subsection{No Posterior Reconstruction in the Main Algorithm}

The main prediction path uses $\bar{\va}_t^-$, not $\bar{\va}_t^+$. This
is because $\bar{\va}_t^-$ is computed before incorporating $\vy(t)$,
whereas $\bar{\va}_t^+$ already incorporates the centered observation
embedding $\tilde{\vb}_t$. A reconstruction based on
$V_B\bar{\va}_t^+$ is therefore a same-time posterior reconstruction in
the centered observation-embedding space, not a causal one-step-ahead
prediction.

Such a posterior reconstruction may be recorded as a diagnostic quantity
or studied in a separate ablation, but it is not part of the main
prediction path. The main read-out and decoder learning objective is
based on
\begin{align}
V_B\bar{\va}_t^-
\quad\longrightarrow\quad
\widehat{\tilde{\vm}}_{t|t-1}
\quad\longrightarrow\quad
\hat{\vm}_{t|t-1}
\quad\longrightarrow\quad
\hat{\vy}_{t|t-1}.
\label{eq:main_causal_prediction_path}
\end{align}

\subsection{Summary of the Staged Prediction Block}

The causal prediction block proceeds in two stages:
\begin{enumerate}
\item \textbf{Feature-level read-out.}
Train $r_{\xi_B}$ to predict centered encoder features from the
operator prediction $V_B\bar{\va}_t^-$ using
\eqref{eq:feature_prediction_loss}.
\item \textbf{Raw-level prediction.}
Restore the encoder-feature mean and train $r_{\xi_B}$ and $g_\alpha$
to predict raw observations using
\eqref{eq:decoder_prediction_loss_with_feature_anchor}. Initially,
$g_\alpha$ may be held fixed and only $r_{\xi_B}$ trained; after the
feature-level alignment stabilizes, $r_{\xi_B}$ and $g_\alpha$ may be
refined together.
\end{enumerate}
The encoder $u_\eta$ and dictionary $\psi_\omega$ remain fixed in the
main formulation. Updating them changes the observation embeddings and
therefore requires a separate controlled refinement stage with
recomputation of the centered variables and periodic re-estimation of
$A^-$, $A^+$, $V_A$, and $V_B$.

\section{Staged Training Summary and Optional Refinement}
\label{sec:staged_training_summary}

The proposed method is trained in stages. This staged design is important
because the embedded-state inference block, the Galerkin operator
estimation block, and the raw-observation decoder block have different
roles. In particular, the projected operators $V_A$ and $V_B$ should be
estimated as Galerkin conditional-operator regressions rather than
absorbed into an unconstrained end-to-end neural predictor.

The main algorithm consists of Phases 0--3 below. An additional
controlled refinement phase may be used after the main algorithm has
converged, but it is not part of the core method.

\subsection{Phase 0: Feature Preparation}

The raw observations are first mapped to encoder features and
observation-embedding coordinates:
\begin{align}
\vm_t
&=
u_\eta(\vy(t)),
&
\vb_t
&=
\psi_\omega(\vm_t).
\label{eq:phase0_features}
\end{align}
The encoder $u_\eta$ and decoder $g_\alpha$ may be initialized by a
standard autoencoder loss,
\begin{align}
\mathcal L_{\mathrm{AE}}(\eta,\alpha)
=
\frac1T
\sum_{t=1}^T
d_{\mathcal Y}
\left(
\vy(t),
g_\alpha(u_\eta(\vy(t)))
\right),
\label{eq:phase0_autoencoder_loss}
\end{align}
although this pretraining is not part of the embedding free-energy
objective.

After the features are computed, store the training means
\begin{align}
\bar{\vm}
&=
\frac1T\sum_{t=1}^T\vm_t,
&
\bar{\vb}
&=
\frac1T\sum_{t=1}^T\vb_t,
\label{eq:phase0_training_means}
\end{align}
and form centered variables
\begin{align}
\tilde{\vm}_t
&=
\vm_t-\bar{\vm},
&
\tilde{\vb}_t
&=
\vb_t-\bar{\vb}.
\label{eq:phase0_centered_features}
\end{align}
The matrix
\begin{align}
\tilde B
=
[\tilde{\vb}_1,\dots,\tilde{\vb}_T]
\label{eq:phase0_centered_B}
\end{align}
is the fixed observation-embedding target used in the core learning
stage. The encoder $u_\eta$ and dictionary $\psi_\omega$ are held fixed
during that stage.

\subsection{Phase 1: Embedded-State and Operator Learning}

Given the fixed centered observation-embedding matrix $\tilde B$, the
prior and posterior embedded-state coordinates $A^-$ and $A^+$ are
estimated by minimizing the finite-dimensional embedding free energy
\eqref{eq:finite_embedding_free_energy}. The projected transfer and
observable operators are updated by the Galerkin ridge-regression
subproblems \eqref{eq:VA_operator_subproblem_revised} and
\eqref{eq:VB_operator_subproblem_revised}.

This phase is the core of the proposed method. It replaces the
CCA-based state construction used in DSE by a variational
embedded-state coordinate learning problem. No inference encoder and no
DSE CCA state warm start are used. The state coordinates are learned
directly as relaxed Galerkin coordinates with intended prior and
posterior embedded-belief semantics.

The output of Phase 1 is
\begin{align}
A^-,
\qquad
A^+,
\qquad
V_A,
\qquad
V_B,
\qquad
R_B,
\qquad
Q_A,
\qquad
\{P_t^-\}_{t=1}^T,
\qquad
P_0,
\qquad
\va_{\mathrm{init}},
\qquad
\bar{\vb},
\qquad
\bar{\vm}.
\label{eq:phase1_outputs}
\end{align}
In the main fixed-metric setting, the metrics are scalar matrices held
fixed during the core optimization. Residual-calibrated and
metric-recursive variants may be considered as extensions of the metric
geometry.

\subsection{Phase 2: Causal Filtering Path Construction}

The training-time coordinates $A^-$ and $A^+$ are variational variables.
They are not themselves used as the main prediction path. After the
operators have been learned, we recompute a causal filtering path using
the same recursion that will be used at prediction time.

The path is initialized by
\begin{align}
\bar{\va}_1^-
=
\va_{\mathrm{init}}.
\label{eq:phase2_initial_prior}
\end{align}
For each time index, use the centered observation embedding
$\tilde{\vb}_t=\vb_t-\bar{\vb}$ and define
\begin{align}
\bar{\va}_t^+
&=
\bar{\va}_t^-
+
K_t
\left(
\tilde{\vb}_t - V_B\bar{\va}_t^-
\right),
\label{eq:phase2_posterior}
\\
\bar{\va}_{t+1}^-
&=
V_A\bar{\va}_t^+.
\label{eq:phase2_prior}
\end{align}
Here $K_t$ is the Galerkin correction gain defined in
\eqref{eq:test_gain}. In the main fixed-metric setting, $P_t^-=pI$ is
held fixed. In a metric-recursive variant, $P_t^-$ and $P_t^+$ may be
updated by \eqref{eq:test_cov_update_simple}--\eqref{eq:test_cov_prediction}
or by the Joseph form \eqref{eq:test_cov_update_joseph}.

The resulting sequence
\begin{align}
\bar A^-
&=
[\bar{\va}_1^-,\dots,\bar{\va}_T^-],
&
\bar A^+
&=
[\bar{\va}_1^+,\dots,\bar{\va}_T^+]
\label{eq:phase2_filtered_paths}
\end{align}
is the causal filtered coordinate sequence associated with the learned
operators. It is distinct from the training-time variational coordinates
$A^-$ and $A^+$ and is the path used for prediction and read-out
training.

\subsection{Phase 3: Causal Read-Out and Decoder Learning}

The main prediction path uses the prior coordinate $\bar{\va}_t^-$,
which is computed before incorporating $\vy(t)$. Under the centered
observation convention, the predicted observation embedding is
\begin{align}
\widehat{\tilde{\vb}}_{t|t-1}
=
V_B\bar{\va}_t^-.
\label{eq:phase3_predicted_embedding}
\end{align}
A read-out map $r_{\xi_B}$ maps this operator-projected embedding to a
centered encoder-feature prediction:
\begin{align}
\widehat{\tilde{\vm}}_{t|t-1}
=
r_{\xi_B}(V_B\bar{\va}_t^-).
\label{eq:phase3_predicted_feature}
\end{align}
The feature-level causal prediction loss is
\begin{align}
\mathcal L_{\mathrm{feat}}^{\mathrm{pred}}(\xi_B)
=
\frac1T
\sum_{t=1}^T
\left\|
\tilde{\vm}_t
-
r_{\xi_B}(V_B\bar{\va}_t^-)
\right\|_2^2
+
\lambda_{r,B}\Omega(\xi_B).
\label{eq:phase3_feature_loss}
\end{align}
This loss aligns the operator-predicted centered observation embedding
with the centered encoder-feature space.

Before decoding, the encoder-feature mean is restored:
\begin{align}
\hat{\vm}_{t|t-1}
=
\bar{\vm}
+
\widehat{\tilde{\vm}}_{t|t-1}
=
\bar{\vm}
+
r_{\xi_B}(V_B\bar{\va}_t^-).
\label{eq:phase3_mean_restored_feature}
\end{align}
The decoder then maps the predicted feature to the raw observation
space:
\begin{align}
\hat{\vy}_{t|t-1}
=
g_\alpha\!\left(
\bar{\vm}
+
r_{\xi_B}(V_B\bar{\va}_t^-)
\right).
\label{eq:phase3_raw_prediction}
\end{align}
The raw causal prediction loss is
\begin{align}
\mathcal L_{\mathrm{raw}}^{\mathrm{pred}}(\xi_B,\alpha)
=
\frac1T
\sum_{t=1}^T
d_{\mathcal Y}
\left(
\vy(t),
g_\alpha\!\left(
\bar{\vm}
+
r_{\xi_B}(V_B\bar{\va}_t^-)
\right)
\right)
+
\lambda_g\Omega(\alpha)
+
\lambda_{r,B}\Omega(\xi_B).
\label{eq:phase3_raw_loss}
\end{align}
In practice we retain the centered feature-level anchor and use
\begin{align}
\mathcal L_{\mathrm{readout}}
=
\mathcal L_{\mathrm{raw}}^{\mathrm{pred}}
+
\lambda_m
\mathcal L_{\mathrm{feat}}^{\mathrm{pred}}.
\label{eq:phase3_readout_loss}
\end{align}
The feature anchor prevents the decoder from absorbing arbitrary changes
in the read-out representation.

The main algorithm does not include posterior reconstruction from
$V_B\bar{\va}_t^+$. Since $\bar{\va}_t^+$ already incorporates the
centered observation embedding $\tilde{\vb}_t$, reconstruction from
$V_B\bar{\va}_t^+$ is a same-time posterior reconstruction rather than a
causal prediction. Such a quantity may be recorded for diagnostics or
studied in a separate ablation, but it is not part of the main
formulation.

\subsection{Main Algorithm Summary}

The main method is summarized in
Algorithm~\ref{alg:staged_embedding_free_energy}. It consists of
feature preparation, core embedded-state and operator learning, causal
filtering path construction, and causal read-out / decoder learning.
The optional refinement described in the next subsection is not required
for the method to be well-defined.

\begin{algorithm}[t]
\caption{Staged Embedding-Free-Energy Operator Learning}
\label{alg:staged_embedding_free_energy}
\begin{algorithmic}[1]
\Require Observations $\{\vy(t)\}_{t=1}^T$, encoder $u_\eta$,
observation dictionary $\psi_\omega$, decoder $g_\alpha$, dimensions
$d_A,d_B$, scalar metric hyperparameters, ridge parameters.
\Ensure Operators $V_A,V_B$, causal filtering rule, read-out
$r_{\xi_B}$, decoder $g_\alpha$, stored means $\bar{\vb},\bar{\vm}$,
and prediction map.

\State \textbf{Phase 0: Feature preparation.}
\State Compute $\vm_t\gets u_\eta(\vy(t))$ and
$\vb_t\gets \psi_\omega(\vm_t)$.
\State Compute and store training means $\bar{\vm}$ and $\bar{\vb}$.
\State Form centered variables
$\tilde{\vm}_t\gets\vm_t-\bar{\vm}$ and
$\tilde{\vb}_t\gets\vb_t-\bar{\vb}$.
\State Form $\tilde B\gets[\tilde{\vb}_1,\dots,\tilde{\vb}_T]$.

\State \textbf{Phase 1: Core embedded-state / operator learning.}
\State Provide nondegenerate numerical seeds for $A^+$ and $A^-$;
these seeds are implementation choices and do not define the
embedded-state coordinates.
\State Set the initial prior coordinate $\va_{\mathrm{init}}$; if no
meaningful initial embedded law is available, use the origin of the
centered Galerkin gauge.
\State Minimize the embedding free energy by alternating coordinate
updates for $A^+$ and $A^-$ with closed-form Galerkin ridge updates for
$V_A$ and $V_B$.
\State Use $\tilde B$ in the observable-operator update.
\State Optionally calibrate $R_B$ and $Q_A$ from residuals in an outer
loop.

\State \textbf{Phase 2: Causal filtering path.}
\State Initialize $\bar{\va}_1^-\gets\va_{\mathrm{init}}$.
\State Run the learned filtering recursion with centered innovations
$\tilde{\vb}_t-V_B\bar{\va}_t^-$ to compute $\bar A^-$ and $\bar A^+$.

\State \textbf{Phase 3: Read-out / decoder learning.}
\State Train $r_{\xi_B}$ to predict centered encoder features
$\tilde{\vm}_t$ using the feature-level causal prediction loss.
\State Restore the encoder-feature mean and train $r_{\xi_B}$ and
$g_\alpha$ using the raw causal prediction loss with feature anchor.

\State \Return $V_A$, $V_B$, $r_{\xi_B}$, $g_\alpha$,
$\bar{\vb}$, $\bar{\vm}$, and the causal filtering / prediction rule.
\end{algorithmic}
\end{algorithm}

\subsection{Optional Controlled Refinement}

The main formulation keeps $u_\eta$ and $\psi_\omega$ fixed during the
embedded-state and operator-learning block. This preserves the meaning
of the centered observation-embedding matrix $\tilde B$ and avoids
co-adaptation in which the encoder absorbs temporal structure and the
learned operators become degenerate.

After the main staged procedure has converged, one may optionally refine
the feature maps and decoder in a controlled outer loop. Such a
refinement should not be performed as fully joint training from scratch.
If $\psi_\omega$ or $u_\eta$ is updated, the observation embeddings
$\vb_t=\psi_\omega(u_\eta(\vy(t)))$ change, and therefore the training
means, centered variables, embedded-state coordinates, and operators
must be recomputed or updated. A controlled refinement loop has the
following form:
\begin{enumerate}
\item update selected neural modules, such as $r_{\xi_B}$, $g_\alpha$,
$\psi_\omega$, or $u_\eta$, for a small number of gradient steps using
the causal prediction loss;
\item recompute $\bar{\vb}$, $\bar{\vm}$, $\tilde B$, and $\tilde M$ if
$u_\eta$ or $\psi_\omega$ has changed;
\item rerun a small number of core learning iterations to update
$A^-$, $A^+$, $V_A$, and $V_B$;
\item recompute the causal filtering path $\bar A^-$ and $\bar A^+$.
\end{enumerate}
This optional stage is a performance refinement, not part of the
mathematical core of the method.

% Acknowledgements and Disclosure of Funding should go at the end, before appendices and references

\acks{

}

% Manual newpage inserted to improve layout of sample file - not
% needed in general before appendices/bibliography.

\newpage


\appendix

\section{Derivation of the Hilbert-Space Correction Update}
\label{app:hilbert_correction_derivation}

This appendix gives a more detailed derivation of the posterior
correction objective used in Section~\ref{sec:statistical_objective}.
The goal is to clarify that the update is a Hilbert-space
residual-minimization principle for embedded beliefs, rather than a
Gaussian likelihood assumption on finite-dimensional latent
coordinates.

\subsection{Conditional Mean Embeddings as Hilbert-Space Projections}

Let
\begin{align}
U_t
:=
\phi_x(\rvx(t))
\in
L^2(\Omega;\mathcal H)
\end{align}
be the Hilbert-space valued latent feature at time $t$. Let
\begin{align}
\mathcal F_t
:=
\sigma(\rvy(1),\dots,\rvy(t))
\end{align}
be the observation sigma-algebra up to time $t$. The posterior embedded
state law is the Hilbert-space valued conditional expectation
\begin{align}
\mu_{x,t}^{+,\star}
=
\mathbb E[
U_t
\mid
\mathcal F_t
].
\end{align}
It is characterized as the orthogonal projection of $U_t$ onto the
closed subspace of $\mathcal F_t$-measurable Hilbert-space valued random
variables:
\begin{align}
\mu_{x,t}^{+,\star}
=
\arg\min_{M\in L^2(\mathcal F_t;\mathcal H)}
\mathbb E
\left[
\|U_t-M\|_{\mathcal H}^2
\right].
\label{eq:app_conditional_expectation_projection}
\end{align}
Similarly, the prior embedded state law is
\begin{align}
\mu_{x,t}^{-,\star}
=
\mathbb E[
U_t
\mid
\mathcal F_{t-1}
].
\end{align}
This exact characterization is not directly computable in our setting,
because the latent state $\rvx(t)$ is not observed and the full
conditional expectation space is infinite-dimensional. The proposed
method therefore estimates the projected embedded beliefs in a finite
Galerkin trial space by using observable embedding consistency and
transfer consistency.

\subsection{Observation Consistency}

At time $t$, the observation $\vy(t)$ gives the empirical observation-law
embedding
\begin{align}
\psi_y(\vy(t))
\in \mathcal G.
\end{align}
Its projected Galerkin representative is
\begin{align}
\psi_t^\Pi
=
P_B\psi_y(\vy(t))
\in \mathcal G_B.
\end{align}
Let
\begin{align}
\mathcal O_B:\mathcal H_A\to\mathcal G_B
\end{align}
be the projected observable operator. If $\mu\in\mathcal H_A$ is a
candidate posterior embedded state, then $\mathcal O_B\mu$ is the
predicted projected observation-law embedding associated with that
candidate state. Therefore, a natural observation-consistency condition
is
\begin{align}
\mathcal O_B\mu
\approx
\psi_t^\Pi.
\label{eq:app_observation_consistency}
\end{align}

In an ideal noiseless and exactly represented model, one could impose
the hard constraint
\begin{align}
\mathcal O_B\mu
=
\psi_t^\Pi.
\label{eq:app_hard_observation_constraint}
\end{align}
Among all elements satisfying this constraint, the corrected embedded
state should be as close as possible to the prior embedded state
$\mu_t^-$. This gives the constrained projection problem
\begin{align}
\min_{\mu\in\mathcal H_A}
\frac12
\|\mu-\mu_t^-\|_{(\mathcal C_t^-)^\dagger}^2
\quad
\text{subject to}
\quad
\mathcal O_B\mu=\psi_t^\Pi.
\label{eq:app_hard_projection_problem}
\end{align}
Here $\mathcal C_t^-$ is a positive semi-definite metric or
preconditioning operator associated with the prior embedded state. It
specifies the local geometry of the correction step: directions with
larger uncertainty are allowed to move more, while directions with
smaller uncertainty are more strongly constrained.

The hard constraint in \eqref{eq:app_hard_projection_problem} is too
strong in finite-dimensional learned feature spaces. The observation is
stochastic, the Galerkin projection is approximate, and the observable
operator is estimated from data. We therefore relax
\eqref{eq:app_hard_projection_problem} into the Tikhonov-regularized
problem
\begin{align}
\mu_t^+
=
\arg\min_{\mu\in\mathcal H_A}
\frac12
\|\psi_t^\Pi-\mathcal O_B\mu\|_{\mathcal R_B^{-1}}^2
+
\frac12
\|\mu-\mu_t^-\|_{(\mathcal C_t^-)^\dagger}^2
+
\Omega_A(\mu).
\label{eq:app_tikhonov_correction_objective}
\end{align}
The first term penalizes violation of the observation-consistency
condition. The second term is a prior-to-posterior proximal penalty in
the embedded-state space. The optional term $\Omega_A$ may be used for
additional coordinate regularization. In the main formulation,
$\Omega_A=0$.

The metrics $\mathcal R_B$ and $\mathcal C_t^-$ should not be interpreted
as requiring a Gaussian latent state-space model. They specify the
geometry of the embedding residual and the correction step. Thus
\eqref{eq:app_tikhonov_correction_objective} is a loss-based
Hilbert-space update for embedded beliefs, not a negative log density of
an explicit parametric latent model.

\subsection{Normal Equation and Innovation Form}

Assume $\Omega_A=0$. The objective in
\eqref{eq:app_tikhonov_correction_objective} becomes
\begin{align}
J_t(\mu)
=
\frac12
\|\psi_t^\Pi-\mathcal O_B\mu\|_{\mathcal R_B^{-1}}^2
+
\frac12
\|\mu-\mu_t^-\|_{(\mathcal C_t^-)^\dagger}^2.
\label{eq:app_quadratic_correction_objective}
\end{align}
Taking the first variation gives
\begin{align}
-
\mathcal O_B^\ast
\mathcal R_B^{-1}
\left(
\psi_t^\Pi-\mathcal O_B\mu
\right)
+
(\mathcal C_t^-)^\dagger
(\mu-\mu_t^-)
=
0.
\label{eq:app_first_variation}
\end{align}
Thus the minimizer satisfies the normal equation
\begin{align}
\left(
\mathcal O_B^\ast
\mathcal R_B^{-1}
\mathcal O_B
+
(\mathcal C_t^-)^\dagger
\right)
\mu_t^+
=
\mathcal O_B^\ast
\mathcal R_B^{-1}
\psi_t^\Pi
+
(\mathcal C_t^-)^\dagger
\mu_t^-.
\label{eq:app_normal_equation_hilbert}
\end{align}
Equivalently, on the subspaces where the relevant inverses are
well-defined, the solution can be written in innovation form:
\begin{align}
\mu_t^+
&=
\mu_t^-
+
\mathcal K_t
\left(
\psi_t^\Pi-\mathcal O_B\mu_t^-
\right),
\label{eq:app_innovation_form_hilbert}
\\
\mathcal K_t
&=
\mathcal C_t^-
\mathcal O_B^\ast
\left(
\mathcal O_B\mathcal C_t^-\mathcal O_B^\ast
+
\mathcal R_B
\right)^{-1}.
\label{eq:app_gain_hilbert}
\end{align}
This follows by applying the Woodbury identity to
\eqref{eq:app_normal_equation_hilbert}. The residual
\begin{align}
\psi_t^\Pi-\mathcal O_B\mu_t^-
\end{align}
is the projected observation innovation, and $\mathcal K_t$ is the
Hilbert-space correction gain.

This is the precise sense in which the update is KKR-like: it has a
prior embedded belief, an observation innovation, and a gain determined
by a residual metric and a prior update metric. However, the update is
not presented as an exact kernel Bayes rule and is not derived from a
Gaussian density over the Galerkin coordinates.

\subsection{Finite-Dimensional Coordinate Form}

Let
\begin{align}
\mu_t^-
&=
\sum_{i=1}^{d_A} a_{i,t}^- \chi_i,
&
\mu_t^+
&=
\sum_{i=1}^{d_A} a_{i,t}^+ \chi_i.
\end{align}
On the observation side, the main formulation uses the centered
Galerkin convention. Thus, if
\begin{align}
\psi_t^\Pi
=
\sum_{j=1}^{d_B} b_{j,t}\omega_j
\end{align}
is the projected observation embedding, its centered coordinate is
denoted by
\begin{align}
\tilde{\vb}_t
=
\vb_t-\bar{\vb}.
\end{align}
The finite-dimensional correction objective corresponding to
\eqref{eq:app_tikhonov_correction_objective} is
\begin{align}
\va_t^+
=
\arg\min_{\va\in\mathbb R^{d_A}}
\frac12
\|\tilde{\vb}_t-V_B\va\|_{R_B^{-1}}^2
+
\frac12
\|\va-\va_t^-\|_{(P_t^-)^\dagger}^2
+
\Omega_A(\va).
\label{eq:app_finite_correction_objective}
\end{align}
When $\Omega_A=0$, the finite-dimensional normal equation is
\begin{align}
\left(
V_B^\top R_B^{-1}V_B
+
(P_t^-)^\dagger
\right)
\va_t^+
=
V_B^\top R_B^{-1}\tilde{\vb}_t
+
(P_t^-)^\dagger\va_t^-.
\label{eq:app_finite_normal_equation}
\end{align}
The corresponding innovation form is
\begin{align}
\va_t^+
&=
\va_t^-
+
K_t
\left(
\tilde{\vb}_t-V_B\va_t^-
\right),
\label{eq:app_finite_innovation}
\\
K_t
&=
P_t^-V_B^\top
\left(
V_BP_t^-V_B^\top+R_B
\right)^{-1}.
\label{eq:app_finite_gain}
\end{align}
This is the correction step used in the causal filtering recursion. The
use of $\tilde{\vb}_t$ reflects the centered Galerkin coordinate
convention; it is not an additional probabilistic assumption on the
latent process.


\section{Coordinate-Update Derivations}
\label{app:coordinate_update_derivations}

This appendix derives the coordinate and operator updates used in
Algorithm~\ref{alg:core_embedding_learning}. Throughout this appendix we
use the finite-dimensional embedding free energy
\eqref{eq:finite_embedding_free_energy}. We focus on the main case
$\Omega_A=0$. If a nonzero convex regularizer $\Omega_A$ is used, the
posterior-coordinate update becomes the corresponding proximal or
constrained quadratic subproblem.

For compact notation, define
\begin{align}
M_t
:=
(P_t^-)^\dagger,
\qquad
S_B
:=
R_B^{-1},
\qquad
S_A
:=
Q_A^{-1}.
\label{eq:app_metric_short_notation}
\end{align}

\subsection{Posterior-Coordinate Update}

Fix $A^-$, $V_A$, and $V_B$. The terms in
\eqref{eq:finite_embedding_free_energy} that depend on a single
posterior coordinate $\va_t^+$ are
\begin{align}
J_t^+(\va)
&=
\frac12
\|\tilde{\vb}_t - V_B\va\|_{R_B^{-1}}^2
+
\frac12
\|\va-\va_t^-\|_{(P_t^-)^\dagger}^2
\nonumber\\
&\quad+
\mathbf 1_{\{t<T\}}
\frac12
\|\va_{t+1}^- - V_A\va\|_{Q_A^{-1}}^2 .
\label{eq:app_Aplus_local_objective}
\end{align}
Using the shorthand in \eqref{eq:app_metric_short_notation}, this is
\begin{align}
J_t^+(\va)
&=
\frac12
(\tilde{\vb}_t - V_B\va)^\top
S_B
(\tilde{\vb}_t - V_B\va)
+
\frac12
(\va-\va_t^-)^\top
M_t
(\va-\va_t^-)
\nonumber\\
&\quad+
\mathbf 1_{\{t<T\}}
\frac12
(\va_{t+1}^- - V_A\va)^\top
S_A
(\va_{t+1}^- - V_A\va).
\label{eq:app_Aplus_local_objective_expanded}
\end{align}
Taking the gradient with respect to $\va$ gives
\begin{align}
\nabla_\va J_t^+(\va)
&=
V_B^\top S_B
(V_B\va-\tilde{\vb}_t)
+
M_t(\va-\va_t^-)
\nonumber\\
&\quad+
\mathbf 1_{\{t<T\}}
V_A^\top S_A
(V_A\va-\va_{t+1}^-).
\label{eq:app_Aplus_gradient}
\end{align}
Setting this gradient to zero yields
\begin{align}
\left[
V_B^\top S_B V_B
+
M_t
+
\mathbf 1_{\{t<T\}}
V_A^\top S_A V_A
\right]\va_t^+
=
V_B^\top S_B\tilde{\vb}_t
+
M_t\va_t^-
+
\mathbf 1_{\{t<T\}}
V_A^\top S_A\va_{t+1}^-.
\label{eq:app_Aplus_normal_equation_short}
\end{align}
Therefore,
\begin{align}
H_t^+
&=
V_B^\top R_B^{-1}V_B
+
(P_t^-)^\dagger
+
\mathbf 1_{\{t<T\}}
V_A^\top Q_A^{-1}V_A,
\label{eq:app_Aplus_H}
\\
h_t^+
&=
V_B^\top R_B^{-1}\tilde{\vb}_t
+
(P_t^-)^\dagger\va_t^-
+
\mathbf 1_{\{t<T\}}
V_A^\top Q_A^{-1}\va_{t+1}^-,
\label{eq:app_Aplus_h}
\end{align}
and the posterior-coordinate update is
\begin{align}
\va_t^+
=
(H_t^+)^{-1}h_t^+.
\label{eq:app_Aplus_solution}
\end{align}
When $H_t^+$ is ill-conditioned, the system can be solved with an
additional small diagonal loading or by using a numerically stable
linear solver rather than explicitly forming the inverse.

\subsection{Prior-Coordinate Update}

Next fix $A^+$ and $V_A$. For $t>1$, the terms depending on
$\va_t^-$ are
\begin{align}
J_t^-(\va)
&=
\frac12
\|\va_t^+ - \va\|_{(P_t^-)^\dagger}^2
+
\frac12
\|\va - V_A\va_{t-1}^+\|_{Q_A^{-1}}^2 .
\label{eq:app_Aminus_local_objective}
\end{align}
Equivalently,
\begin{align}
J_t^-(\va)
&=
\frac12
(\va_t^+ - \va)^\top
M_t
(\va_t^+ - \va)
+
\frac12
(\va - V_A\va_{t-1}^+)^\top
S_A
(\va - V_A\va_{t-1}^+).
\label{eq:app_Aminus_local_objective_expanded}
\end{align}
Taking the gradient with respect to $\va$ gives
\begin{align}
\nabla_\va J_t^-(\va)
=
M_t(\va-\va_t^+)
+
S_A(\va - V_A\va_{t-1}^+).
\label{eq:app_Aminus_gradient}
\end{align}
Setting the gradient to zero gives
\begin{align}
(M_t+S_A)\va_t^-
=
M_t\va_t^+
+
S_A V_A\va_{t-1}^+.
\label{eq:app_Aminus_normal_equation_short}
\end{align}
Thus,
\begin{align}
H_t^-
&=
(P_t^-)^\dagger
+
Q_A^{-1},
\label{eq:app_Aminus_H}
\\
h_t^-
&=
(P_t^-)^\dagger\va_t^+
+
Q_A^{-1}V_A\va_{t-1}^+,
\label{eq:app_Aminus_h}
\end{align}
and
\begin{align}
\va_t^-
=
(H_t^-)^{-1}h_t^-,
\qquad
t=2,\dots,T.
\label{eq:app_Aminus_solution}
\end{align}

For $t=1$, the transfer-consistency term is replaced by the initial
prior term. The local objective becomes
\begin{align}
J_1^-(\va)
&=
\frac12
\|\va_1^+ - \va\|_{(P_1^-)^\dagger}^2
+
\frac12
\|\va - \va_{\mathrm{init}}\|_{P_0^{-1}}^2.
\label{eq:app_Aminus_initial_objective}
\end{align}
The corresponding normal equation is
\begin{align}
\left(
(P_1^-)^\dagger + P_0^{-1}
\right)\va_1^-
=
(P_1^-)^\dagger\va_1^+
+
P_0^{-1}\va_{\mathrm{init}}.
\label{eq:app_Aminus_initial_normal}
\end{align}
Hence
\begin{align}
H_1^-
&=
(P_1^-)^\dagger+P_0^{-1},
\label{eq:app_Aminus_initial_H}
\\
h_1^-
&=
(P_1^-)^\dagger\va_1^+
+
P_0^{-1}\va_{\mathrm{init}},
\label{eq:app_Aminus_initial_h}
\\
\va_1^-
&=
(H_1^-)^{-1}h_1^-.
\label{eq:app_Aminus_initial_solution}
\end{align}

For data consisting of multiple trajectories, the formula
\eqref{eq:app_Aminus_initial_solution} is applied at the beginning of
each trajectory. Transfer-consistency terms are formed only within each
trajectory and never across trajectory boundaries.

\subsection{Transfer-Operator Update}

For fixed $A^-$ and $A^+$, the transfer matrix $V_A$ is estimated from
the soft relation
\begin{align}
\va_{t+1}^-
\approx
V_A\va_t^+.
\end{align}
Let
\begin{align}
X_A
&:=
A_{1:T-1}^+
=
[\va_1^+,\dots,\va_{T-1}^+],
&
Y_A
&:=
A_{2:T}^-
=
[\va_2^-,\dots,\va_T^-].
\label{eq:app_transfer_design_matrices}
\end{align}
For multiple trajectories, $X_A$ and $Y_A$ are formed by concatenating
only valid within-trajectory transition pairs. No column pair is formed
from the last time index of one trajectory to the first time index of
another trajectory.

The default operator subproblem is the ridge regression
\begin{align}
V_A
=
\arg\min_{V\in\mathbb R^{d_A\times d_A}}
\frac12
\|Y_A - VX_A\|_F^2
+
\frac{\lambda_A}{2}
\|V\|_F^2.
\label{eq:app_transfer_ridge_problem}
\end{align}
Taking the derivative with respect to $V$ gives
\begin{align}
(VX_A-Y_A)X_A^\top+\lambda_A V=0.
\label{eq:app_transfer_gradient}
\end{align}
Therefore,
\begin{align}
V(X_AX_A^\top+\lambda_A I)
=
Y_AX_A^\top.
\label{eq:app_transfer_normal_equation}
\end{align}
Assuming $X_AX_A^\top+\lambda_A I$ is invertible, the solution is
\begin{align}
V_A
=
Y_AX_A^\top
(X_AX_A^\top+\lambda_A I)^{-1}.
\label{eq:app_transfer_closed_form}
\end{align}
Substituting the definitions of $X_A$ and $Y_A$ gives
\begin{align}
V_A
=
A_{2:T}^-
(A_{1:T-1}^+)^\top
\left(
A_{1:T-1}^+
(A_{1:T-1}^+)^\top
+
\lambda_A I
\right)^{-1}.
\label{eq:app_VA_closed_form_final}
\end{align}

If a non-isotropic transfer metric $Q_A$ is used directly in the
operator subproblem, the resulting weighted least-squares problem may
lead to a Sylvester-type normal equation unless the coordinates are
whitened or a metric-compatible regularizer is used. The default
algorithm therefore uses scalar, diagonal, or whitened metrics for the
operator updates so that the Galerkin ridge form
\eqref{eq:app_VA_closed_form_final} is preserved.

\subsection{Observable-Operator Update}

For fixed $A^+$, the observable matrix $V_B$ is estimated from the
centered relation
\begin{align}
\tilde{\vb}_t
\approx
V_B\va_t^+.
\end{align}
Let
\begin{align}
X_B
&:=
A^+
=
[\va_1^+,\dots,\va_T^+],
&
Y_B
&:=
\tilde B
=
[\tilde{\vb}_1,\dots,\tilde{\vb}_T].
\label{eq:app_observable_design_matrices}
\end{align}
The observable-operator subproblem is
\begin{align}
V_B
=
\arg\min_{V\in\mathbb R^{d_B\times d_A}}
\frac12
\|Y_B - VX_B\|_F^2
+
\frac{\lambda_B}{2}
\|V\|_F^2.
\label{eq:app_observable_ridge_problem}
\end{align}
The derivative with respect to $V$ is
\begin{align}
(VX_B-Y_B)X_B^\top+\lambda_B V=0.
\label{eq:app_observable_gradient}
\end{align}
Thus,
\begin{align}
V(X_BX_B^\top+\lambda_B I)
=
Y_BX_B^\top,
\label{eq:app_observable_normal_equation}
\end{align}
and
\begin{align}
V_B
=
Y_BX_B^\top
(X_BX_B^\top+\lambda_B I)^{-1}.
\label{eq:app_observable_closed_form}
\end{align}
Equivalently,
\begin{align}
V_B
=
\tilde B(A^+)^\top
\left(
A^+(A^+)^\top
+
\lambda_B I
\right)^{-1}.
\label{eq:app_VB_closed_form_final}
\end{align}

As with the transfer operator, full non-isotropic observation metrics can
be handled by whitening or by solving the corresponding weighted
least-squares system directly. The main formulation keeps the
closed-form Galerkin ridge structure in
\eqref{eq:app_VB_closed_form_final}.

\subsection{Why the Training Coordinates Are Not Hard-Propagated}

During training, the prior coordinates $A^-$ are not defined by the hard
constraint
\begin{align}
\va_{t+1}^-
=
V_A\va_t^+.
\end{align}
Instead, $A^-$ and $A^+$ are optimized as variational embedded-state
coordinates and the transfer relation is imposed softly through the
transfer-consistency energy. This distinction is necessary for the
operator update to be meaningful. If $\va_{t+1}^-$ were defined by
$\va_{t+1}^-=V_A\va_t^+$ during training, the transfer-operator
regression would become tautological.

At test time, in contrast, the learned operator is used causally:
\begin{align}
\bar{\va}_{t+1}^-
=
V_A\bar{\va}_t^+.
\end{align}
Thus, training uses variational coordinate learning, while test-time
prediction uses causal filtering recursion.

\section{Metric Choices and Residual Calibration}
\label{app:metric_choices}

This appendix summarizes the metric choices used in the
finite-dimensional embedding free energy and in the causal filtering
recursion. The matrices
\begin{align}
R_B,\qquad Q_A,\qquad P_t^-,
\qquad P_0
\end{align}
specify the geometry of centered observation residuals, transfer
residuals, prior-to-posterior corrections, and the initial boundary
anchor. They should not be interpreted as covariance parameters of a
linear-Gaussian latent state-space model. Rather, they define weighted
embedding discrepancies and recursive least-squares preconditioners for
Galerkin embedded-belief updates.

\subsection{Roles of the Metric Matrices}

The centered observation residual is
\begin{align}
\rho_t
=
\tilde{\vb}_t - V_B\va_t^+.
\label{eq:app_obs_residual}
\end{align}
The metric $R_B$ weights this residual through
\begin{align}
\|\rho_t\|_{R_B^{-1}}^2.
\end{align}
Thus $R_B$ specifies the geometry of centered observation-embedding
discrepancies in $\mathbb R^{d_B}$.

The transfer residual is
\begin{align}
\epsilon_t
=
\va_{t+1}^- - V_A\va_t^+.
\label{eq:app_transfer_residual}
\end{align}
The metric $Q_A$ weights this residual through
\begin{align}
\|\epsilon_t\|_{Q_A^{-1}}^2.
\end{align}
Thus $Q_A$ specifies the geometry of transfer-consistency errors in the
latent Galerkin coordinate space.

The metric $P_t^-$ weights the prior-to-posterior correction
\begin{align}
\va_t^+ - \va_t^-.
\end{align}
The term
\begin{align}
\|\va_t^+ - \va_t^-\|_{(P_t^-)^\dagger}^2
\end{align}
penalizes deviations from the prior embedded belief. In this sense,
$P_t^-$ plays the role of a local uncertainty metric, inverse-curvature
metric, or recursive least-squares preconditioner for the embedded-state
correction step.

Finally, $P_0$ controls the strength of the initial boundary anchor
\begin{align}
\|\va_1^- - \va_{\mathrm{init}}\|_{P_0^{-1}}^2.
\end{align}
When a meaningful initial embedded law is available, this term anchors
the first prior coordinate to its Galerkin coordinate. When
$\va_{\mathrm{init}}=0$ is used only as a reference prior in the
centered Galerkin gauge, the same term should be interpreted as weak
boundary regularization.

All metric choices affect the geometry of the optimization problem.
They do not change the statistical target, which remains the projected
prior and posterior conditional mean embeddings of the latent state law.

\subsection{Fixed Scalar Metrics}

The main fixed-metric formulation uses scalar metrics
\begin{align}
R_B
&=
\sigma_B^2 I_{d_B},
&
Q_A
&=
\sigma_A^2 I_{d_A},
&
P_t^-
&=
p I_{d_A},
&
P_0
&=
c_0p I_{d_A}.
\label{eq:app_fixed_scalar_metrics}
\end{align}
With this choice, the embedding free energy contains the weighted
Euclidean terms
\begin{align}
\frac{1}{2\sigma_B^2}
\sum_{t=1}^T
\|\tilde{\vb}_t - V_B\va_t^+\|_2^2
+
\frac{1}{2p}
\sum_{t=1}^T
\|\va_t^+ - \va_t^-\|_2^2
+
\frac{1}{2\sigma_A^2}
\sum_{t=1}^{T-1}
\|\va_{t+1}^- - V_A\va_t^+\|_2^2,
\label{eq:app_fixed_metric_weighted_objective}
\end{align}
together with the initial boundary term
\begin{align}
\frac{1}{2c_0p}
\|\va_1^- - \va_{\mathrm{init}}\|_2^2.
\label{eq:app_fixed_metric_initial_term}
\end{align}
Equivalently, the first three scalar parameters define the loss weights
\begin{align}
w_B=\sigma_B^{-2},
\qquad
w_P=p^{-1},
\qquad
w_A=\sigma_A^{-2}.
\end{align}
The scalar $c_0$ controls the relative strength of the initial boundary
anchor. These scalars are user-specified hyperparameters held fixed
during the core coordinate optimization.

For causal filtering in the same fixed-metric setting, the gain is
\begin{align}
K
=
p V_B^\top
\left(
p V_BV_B^\top+\sigma_B^2 I_{d_B}
\right)^{-1},
\label{eq:app_fixed_metric_gain}
\end{align}
and the posterior correction is
\begin{align}
\bar{\va}_t^+
=
\bar{\va}_t^-
+
K
\left(
\tilde{\vb}_t - V_B\bar{\va}_t^-
\right).
\label{eq:app_fixed_metric_test_update}
\end{align}
Since $P_t^-=pI_{d_A}$ is fixed, this setting does not propagate a
time-varying covariance or inverse-curvature metric.

The metric choices described here are independent of the numerical seed
used to start the alternating coordinate solver. Changing the seed may
affect optimization, but it does not alter the definition of the
embedding free-energy objective or the meaning of the metric terms.

\subsection{Residual-Calibrated Metrics}

A residual-calibrated metric variant can estimate $R_B$ and $Q_A$ from
empirical residuals. Given current coordinates and operators, define
\begin{align}
\rho_t
&=
\tilde{\vb}_t - V_B\va_t^+,
&
\epsilon_t
&=
\va_{t+1}^- - V_A\va_t^+.
\label{eq:app_residuals_for_calibration}
\end{align}
The empirical residual covariance estimates are
\begin{align}
\widehat R_B
&=
\frac1T
\sum_{t=1}^T
\rho_t\rho_t^\top,
\label{eq:app_Rhat_full}
\\
\widehat Q_A
&=
\frac1{T-1}
\sum_{t=1}^{T-1}
\epsilon_t\epsilon_t^\top.
\label{eq:app_Qhat_full}
\end{align}
Directly using these matrices can be unstable, especially when the
number of samples is small or the residuals are low rank. A shrinkage
and diagonal-loading form is therefore
\begin{align}
R_B
&=
(1-\alpha_R)\widehat R_B
+
\alpha_R
\frac{\operatorname{tr}(\widehat R_B)}{d_B}
I_{d_B}
+
\varepsilon_R I_{d_B},
\label{eq:app_RB_shrinkage}
\\
Q_A
&=
(1-\alpha_Q)\widehat Q_A
+
\alpha_Q
\frac{\operatorname{tr}(\widehat Q_A)}{d_A}
I_{d_A}
+
\varepsilon_Q I_{d_A}.
\label{eq:app_QA_shrinkage}
\end{align}
A safer first residual-calibrated variant uses only the diagonal parts:
\begin{align}
R_B
&=
\operatorname{diag}(\widehat R_B)
+
\varepsilon_R I_{d_B},
\label{eq:app_RB_diagonal}
\\
Q_A
&=
\operatorname{diag}(\widehat Q_A)
+
\varepsilon_Q I_{d_A}.
\label{eq:app_QA_diagonal}
\end{align}
The residual-calibrated metrics should be updated in an outer loop.
During each coordinate-minimization step, $R_B$ and $Q_A$ are held
fixed. Treating them as unconstrained inner-loop optimization variables
can lead to scale degeneracies and is not part of the main formulation.

\subsection{Interaction with Operator Updates}

The coordinate updates can use the metrics $R_B$ and $Q_A$ directly.
However, if full non-isotropic metrics are also used inside the
operator subproblems, the closed-form ridge updates may change.

For example, the standard centered observable-operator update is
\begin{align}
V_B
=
\tilde B(A^+)^\top
\left(
A^+(A^+)^\top+\lambda_B I
\right)^{-1}.
\label{eq:app_standard_centered_VB_update}
\end{align}
If the operator subproblem is instead written with a full observation
metric as
\begin{align}
\min_V
\frac12
\sum_{t=1}^T
\|\tilde{\vb}_t - V\va_t^+\|_{R_B^{-1}}^2
+
\frac{\lambda_B}{2}
\|V\|_F^2,
\label{eq:app_metric_weighted_operator_problem}
\end{align}
then the normal equation becomes
\begin{align}
R_B^{-1}(VA^+ - \tilde B)(A^+)^\top
+
\lambda_B V
=
0.
\label{eq:app_metric_weighted_operator_normal}
\end{align}
After multiplying by $R_B$, this becomes
\begin{align}
V A^+(A^+)^\top
-
\tilde B(A^+)^\top
+
\lambda_B R_B V
=
0,
\label{eq:app_metric_weighted_operator_sylvester}
\end{align}
which is generally not the same ridge form unless the metric is scalar,
diagonal in a compatible coordinate system, or the coordinates are
whitened. For this reason, the main formulation keeps the operator
subproblems in the standard Galerkin ridge form. Full weighted
least-squares operator updates may be studied as an implementation
variant.

\subsection{KKR-Style Metric-Recursive Variant}

A closer KKR-style metric variant propagates $P_t^-$ as a
time-dependent covariance or inverse-curvature metric. Given $P_t^-$,
the gain is
\begin{align}
K_t
=
P_t^-V_B^\top
\left(
V_BP_t^-V_B^\top+R_B
\right)^{-1}.
\label{eq:app_metric_recursive_gain}
\end{align}
The posterior coordinate is updated by
\begin{align}
\bar{\va}_t^+
=
\bar{\va}_t^-
+
K_t
\left(
\tilde{\vb}_t - V_B\bar{\va}_t^-
\right).
\label{eq:app_metric_recursive_coordinate_update}
\end{align}
The corresponding inverse-curvature update can be written as
\begin{align}
P_t^+
=
P_t^-
-
K_tV_BP_t^-.
\label{eq:app_metric_recursive_simple_update}
\end{align}
For numerical stability, the Joseph form may be used:
\begin{align}
P_t^+
=
(I-K_tV_B)P_t^-(I-K_tV_B)^\top
+
K_tR_BK_t^\top.
\label{eq:app_metric_recursive_joseph}
\end{align}
The prior metric is propagated by
\begin{align}
P_{t+1}^-
=
V_AP_t^+V_A^\top
+
Q_A.
\label{eq:app_metric_recursive_prediction}
\end{align}
In this formulation, $P_t^-$ should be interpreted as a local metric,
uncertainty matrix, or inverse-curvature object for the Galerkin
embedded-belief update. It is not introduced by assuming that the
Galerkin coordinate itself is distributed according to a Gaussian law.
The recursion is therefore an optional KKR-style metric choice rather
than a change in the statistical target.

\subsection{Numerical Stability}

When using residual-calibrated or metric-recursive variants, the
following numerical precautions are useful:
\begin{align}
R_B
&\leftarrow
\frac12(R_B+R_B^\top)
+
\varepsilon_R I,
\\
Q_A
&\leftarrow
\frac12(Q_A+Q_A^\top)
+
\varepsilon_Q I,
\\
P_t^-
&\leftarrow
\frac12(P_t^-+(P_t^-)^\top)
+
\varepsilon_P I.
\end{align}
If eigenvalues become too small or negative because of numerical error,
eigenvalue flooring can be applied. In long sequences, one may also
control the spectral norm of $V_A$ or use a diagonal or low-rank
approximation to $P_t^-$.

\subsection{Recommended Metric Usage}

The recommended progression is to begin with the fixed scalar metric
formulation:
\begin{align}
R_B=\sigma_B^2 I,
\qquad
Q_A=\sigma_A^2 I,
\qquad
P_t^-=pI,
\qquad
P_0=c_0pI.
\end{align}
After the core solver is stable, residual-calibrated diagonal metrics
can be studied as the main metric refinement. KKR-style
metric-recursive updates should be treated as a further extension or
ablation that strengthens the connection to kernel Kalman filtering
without changing the embedded-belief target.


\section{Initialization and Normalization}
\label{app:initialization_normalization}

This appendix records the limited implementation conventions needed to
start the core alternating solver. These conventions should not be
confused with the statistical target of the method. The embedded-state
coordinates are defined by the embedding free-energy problem and by the
learned causal prediction--correction recursion, not by the numerical
initializer.

\subsection{Initial Prior Coordinate}

The initial prior coordinate is conceptually distinct from the numerical
initialization of the coordinate arrays $A^-$ and $A^+$. Ideally,
\begin{align}
\va_{\mathrm{init}}
\approx
P_A\mu_{\mathrm{init}},
\label{eq:app_initial_prior_coordinate}
\end{align}
where $\mu_{\mathrm{init}}$ is an initial prior embedded state law. If an
initial ensemble, an externally specified initial distribution, or
another reliable estimate of this embedded law is available, its
Galerkin coordinate may be supplied to the algorithm.

If no meaningful initial embedded law is available, we use
$\va_{\mathrm{init}}=0$ in the centered Galerkin gauge. This zero vector
is a reference prior, not a probabilistic claim about the true initial
state law. In that case, the initial term in the embedding free energy
acts as weak boundary regularization. Its strength is controlled by
$P_0=c_0pI$.

\subsection{Numerical Seeds for the Training Coordinates}

The coordinate arrays
\begin{align}
A^-
&=
[\va_1^-,\dots,\va_T^-],
&
A^+
&=
[\va_1^+,\dots,\va_T^+]
\end{align}
are optimized as training-time variational coordinates. Since the true
latent states are not observed, these arrays cannot be initialized from
samples of $\phi_x(\vx(t))$. They require numerical seeds only because
the embedding free-energy problem is solved by an alternating
optimization procedure.

These seeds are not part of the statistical definition of the state.
They do not define the latent coordinate system, and they should not be
interpreted as prior or posterior embedded beliefs. In particular, the
observation embedding matrix $\tilde B$ is not used to define
$A^-$ or $A^+$. It enters the learning problem through the
observation-consistency and correction terms.

The implementation must provide nondegenerate numerical seeds for
$A^{-,(0)}$ and $A^{+,(0)}$. The exact choice of the seed is an
implementation protocol and may be deterministic, randomized, or
externally supplied. The only requirement at the formulation level is
that the seed should allow the alternating solver to start without
degenerate operator updates.

\subsection{No DSE CCA State Warm Start}

The proposed method replaces the CCA-based state-construction block of
DSE. Therefore, DSE functional CCA coordinates or stochastic-realization
coordinates are not used as warm starts for $A^-$ or $A^+$. This
separation keeps the proposed state-coordinate learning principle
independent of the DSE state-construction principle.

Spectral initializations such as PCA or SVD of the observation embedding
matrix may be considered as numerical warm starts in separate
implementation studies. They are not part of the main formulation and
should not be described as the state-construction principle of the
method.

\subsection{Observation-Side Centering}

The main formulation uses a centered Galerkin convention on the
observation side. After computing
\begin{align}
\vm_t
=
u_\eta(\vy(t)),
\qquad
\vb_t
=
\psi_\omega(\vm_t),
\end{align}
we define
\begin{align}
\bar{\vb}
&=
\frac1T\sum_{t=1}^T \vb_t,
&
\tilde{\vb}_t
&=
\vb_t-\bar{\vb},
&
\tilde B
&=
[\tilde{\vb}_1,\dots,\tilde{\vb}_T],
\label{eq:app_centered_observation_embeddings}
\\
\bar{\vm}
&=
\frac1T\sum_{t=1}^T \vm_t,
&
\tilde{\vm}_t
&=
\vm_t-\bar{\vm},
&
\tilde M
&=
[\tilde{\vm}_1,\dots,\tilde{\vm}_T].
\label{eq:app_centered_encoder_features}
\end{align}
The core solver uses $\tilde B$. The read-out stage predicts centered
encoder features, and $\bar{\vm}$ is restored before applying the
decoder. Test-time observation embeddings are centered using the stored
training mean $\bar{\vb}$.

This centering convention is a numerical coordinate convention for a
linear no-intercept observable operator. It is not an additional
probabilistic assumption on the latent process.

\subsection{Coordinate Gauge}

The Galerkin coordinates have a linear gauge freedom. For any invertible
matrix $S\in\mathbb R^{d_A\times d_A}$,
\begin{align}
\va_t^\pm
&\mapsto
S\va_t^\pm,
&
V_A
&\mapsto
S V_A S^{-1},
&
V_B
&\mapsto
V_B S^{-1}
\end{align}
represents the same linear coordinate system up to a change of basis.
This non-identifiability is standard in latent-coordinate models and
does not invalidate the operator interpretation.

The main formulation does not rely on hard whitening or affine
re-centering of the latent coordinates. If an implementation applies a
linear rescaling of the latent coordinates for numerical conditioning,
the same coordinate convention must be applied consistently to
$A^-$, $A^+$, $\va_{\mathrm{init}}$, and the operators. In practice, it
is often simpler to recompute $V_A$ and $V_B$ by the Galerkin
ridge-regression updates after any such numerical rescaling.



\vskip 0.2in
\bibliography{ref}


\end{document}
